%%
%% This is file `sample-sigconf-authordraft.tex',
%% generated with the docstrip utility.
%%
%% The original source files were:
%%
%% samples.dtx  (with options: `all,proceedings,bibtex,authordraft')
%% 
%% IMPORTANT NOTICE:
%% 
%% For the copyright see the source file.
%% 
%% Any modified versions of this file must be renamed
%% with new filenames distinct from sample-sigconf-authordraft.tex.
%% 
%% For distribution of the original source see the terms
%% for copying and modification in the file samples.dtx.
%% 
%% This generated file may be distributed as long as the
%% original source files, as listed above, are part of the
%% same distribution. (The sources need not necessarily be
%% in the same archive or directory.)
%%
%%
%% Commands for TeXCount
%TC:macro \cite [option:text,text]
%TC:macro \citep [option:text,text]
%TC:macro \citet [option:text,text]
%TC:envir table 0 1
%TC:envir table* 0 1
%TC:envir tabular [ignore] word
%TC:envir displaymath 0 word
%TC:envir math 0 word
%TC:envir comment 0 0
%%
%%
%% The first command in your LaTeX source must be the \documentclass
%% command.
%%
%% For submission and review of your manuscript please change the
%% command to \documentclass[manuscript, screen, review]{acmart}.
%%
%% When submitting camera ready or to TAPS, please change the command
%% to \documentclass[sigconf]{acmart} or whichever template is required
%% for your publication.
%%
%%
\documentclass[manuscript, screen, sigconf]{acmart}

\usepackage{algorithm}
\usepackage{algpseudocode}
\usepackage{dblfloatfix}
\usepackage{gensymb}
\usepackage{xcolor}
\usepackage{soul}
\usepackage{tikz}
\usetikzlibrary{calc,fit,positioning,bending}
\usepackage{footmisc}

% \usepackage{fontspec}

% \newfontfamily\bodyfont
% [Path=./InterFontFiles/,
% BoldFont=Inter-Bold.ttf,
% ItalicFont=Inter-Italic.ttf,
% BoldItalicFont=Inter-BoldItalic.ttf,
% UprightFont=Inter-Regular.ttf]
% {Inter}

\tikzset{every picture/.style={/utils/exec={\bfseries\footnotesize}}}

\usepackage{enumitem}
%%
%% \BibTeX command to typeset BibTeX logo in the docs
\AtBeginDocument{%
  \providecommand\BibTeX{{%
    Bib\TeX}}}

%% Rights management information.  This information is sent to you
%% when you complete the rights form.  These commands have SAMPLE
%% values in them; it is your responsibility as an author to replace
%% the commands and values with those provided to you when you
%% complete the rights form.
\copyrightyear{2025}
\acmYear{2025}
\setcopyright{cc}
\setcctype{by}
\acmConference[SCF '25]{ACM Symposium on Computational Fabrication}{November 20--21, 2025}{Cambridge, MA, USA}
\acmBooktitle{ACM Symposium on Computational Fabrication (SCF '25), November 20--21, 2025, Cambridge, MA, USA}
\acmDOI{10.1145/3745778.3766660}
\acmISBN{979-8-4007-2034-5/2025/11}
\newcommand{\ergun}[1]{{\color{red} #1}}

%%
%% Submission ID.
%% Use this when submitting an article to a sponsored event. You'll
%% receive a unique submission ID from the organizers
%% of the event, and this ID should be used as the parameter to this command.
%%\acmSubmissionID{123-A56-BU3}

%%
%% For managing citations, it is recommended to use bibliography
%% files in BibTeX format.
%%
%% You can then either use BibTeX with the ACM-Reference-Format style,
%% or BibLaTeX with the acmnumeric or acmauthoryear sytles, that include
%% support for advanced citation of software artefact from the
%% biblatex-software package, also separately available on CTAN.
%%
%% Look at the sample-*-biblatex.tex files for templates showcasing
%% the biblatex styles.
%%

%%
%% The majority of ACM publications use numbered citations and
%% references.  The command \citestyle{authoryear} switches to the
%% "author year" style.
%%
%% If you are preparing content for an event
%% sponsored by ACM SIGGRAPH, you must use the "author year" style of
%% citations and references.
%% Uncommenting
%% the next command will enable that style.
\citestyle{acmauthoryear}


%%
%% end of the preamble, start of the body of the document source.
\begin{document}

%%
%% The "title" command has an optional parameter,
%% allowing the author to define a "short title" to be used in page headers.
\title{Fabricating Interactions through Optical Pathways}

%%
%% The "author" command and its associated commands are used to define
%% the authors and their affiliations.
%% Of note is the shared affiliation of the first two authors, and the
%% "authornote" and "authornotemark" commands
%% used to denote shared contribution to the research.
\author{Himani Deshpande}
% \authornote{Both authors contributed equally to this research.}
\email{hdeshpande11@tamu.edu}
\orcid{0000-0002-4839-549X}
\affiliation{%
  \institution{Computer Science and Engineering 
  Texas A\&M University}
  \city{College Station}
  \state{Texas}
  \country{USA}
}

\author{Tolga Yildiz}
% \authornote{Both authors contributed equally to this research.}
\email{tolgayildiz@tamu.edu}
\orcid{0000-0002-4834-9389}
\affiliation{%
  \institution{Computer Science and Engineering Texas A\&M University}
  \city{College Station}
  \state{Texas}
  \country{USA}
}

\author{Madelein Olivan Villegas}
\email{madeleinvillegas@tamu.edu}
\orcid{0009-0001-5411-7402}
\affiliation{%
  \institution{Computer Science and Engineering Texas A\&M University}
  \city{College Station}
  \state{Texas}
  \country{USA}
}

\author{Ergun Akleman}
\email{ergun@tamu.edu}
\orcid{0000-0003-3618-4166}
\affiliation{%
  \institution{
  % Visual Computing and Computational Media Section, 
  School of Performance, Visualization and Fine Arts,
  % Joint with 
  Computer Science and Engineering, Texas A\&M University}
  \city{College Station}
  \state{Texas}
  \country{USA}
}

\author{Jeeeun Kim}
\email{jeeeun.kim@tamu.edu}
\orcid{0000-0002-8915-481X}
\affiliation{%
  \institution{Computer Science and Engineering,  Texas A\&M University}
  \city{College Station}
  \state{Texas}
  \country{USA}
}

%%
%% By default, the full list of authors will be used in the page
%% headers. Often, this list is too long, and will overlap
%% other information printed in the page headers. This command allows
%% the author to define a more concise list
%% of authors' names for this purpose.
\renewcommand{\shortauthors}{Desphande et al.}

\newcommand{\jk}[1]
{{\fontfamily{cmss}\selectfont \bfseries \color{violet} JK: #1 }}
\newcommand{\hd}[1]
{{\fontfamily{cmss}\selectfont \bfseries \color{blue} HD: #1 }}
\newcommand{\ty}[1]
{{\fontfamily{cmss}\selectfont \bfseries \color{green} TY: #1 }}

\newcommand{\revision}[1]
{{\fontfamily{cmss}\selectfont \bfseries \color{red} #1 }}

%%
%% The abstract is a short summary of the work to be presented in the
%% article.
\begin{abstract}





Transparent materials transmit light without significant scattering or absorption due to \textit{total internal reflection}. 
Transparent channels in 3D printed objects follow this theory, functioning similarly to optical fibers by transmitting input light. 
While prior work enabled basic interactions in custom optical sensors like push and display interactions using photopolymers, complex channeling 
remains challenging, particularly in desktop Fused Deposition Modeling (FDM) due to the inherent printing discontinuities. 
We present a framework that enables low-cost desktop fabrication of optical interactive devices using Voronoi-based segmentation of objects for in-place FDM printing of optical channels. It allows uninterrupted light transmission where the embedded channels are aligned in the XY plane for uninterrupted printing.
We further explore the use of FDM optical objects in routing and channeling strategies to support applications in displays, sensing, and embodied interactions.
%that enable better light transmission with fewer imperfections \jk{fewer than what?}.
%\ergun{Suggestion: When I read the paper and this part I did not find a clear list of the techniques. I will start here by saying the following: \textbf{We introduce two techniques to enable the design and fabrication of interactive devices using light signal channeling and routing}: (1)
%\textbf{In-place printing with Voronoi-based segmentation:}
%This technique prints optical channels directly into the object using FDM 3D printing, with segmentation strategies that allow uninterrupted light paths 
%\textbf{(2) Hybrid fabrication:}
%This involves inserting pre-printed transparent optical channels (especially using flexible materials like TPU) into a 3D-printed shell. This method allows better light transmission due to fewer printing imperfections but requires careful design of the housing and insertion paths.}
% This method not only ensures continuity in light routing but also facilitates user-defined interactions through modular design. 
% We employ a Voronoi-based segmentation algorithm to simplify 3D printing of interactive optical devices using desktop FDM.


\end{abstract}

%%
%% The code below is generated by the tool at http://dl.acm.org/ccs.cfm.
%% Please copy and paste the code instead of the example below.
%%
\begin{CCSXML}
<ccs2012>
   <concept>
       <concept_id>10003120.10003121</concept_id>
       <concept_desc>Human-centered computing~Human computer interaction (HCI)</concept_desc>
       <concept_significance>500</concept_significance>
       </concept>
 </ccs2012>
\end{CCSXML}

\ccsdesc[500]{Human-centered computing~Human computer interaction (HCI)}


%%
%% Keywords. The author(s) should pick words that accurately describe
%% the work being presented. Separate the keywords with commas.
\keywords{FDM, Modeling, Optical Channels}
%% A "teaser" image appears between the author and affiliation
%% information and the body of the document, and typically spans the
%% page.
\begin{teaserfigure}
  \centering
\includegraphics[width=\textwidth]{figures/workflow.png}
  \caption{Two curves (a), are used for sampling 5 light channels (b), which are used for Voronoi partitioning (c), each partition can be assigned to a different material and fabrication method (d), optical channels are printed in transparent PLA in-place with the black shell in PLA and blue shell in TPU. Modularity enables material based interactivity with optical elements (e).}
  \Description{A six-part figure (a–f) illustrating a hybrid fabrication workflow for optical elements using Voronoi-based partitioning. (a) Two curves define an input light profile (blue line) and output light profile (red semicircle). (b) Five light-guiding channels are drawn between the input and output profiles (green lines). (c) A Voronoi diagram segments the volume into partitions based on the channel layout. (d) Partitions are separated and prepared for fabrication. Black PLA partitions, and blue TPU partition have transparent PLA optical channels printed in-place; (f) The final assembled object demonstrates altered light behavior, where the same input line that produces a semicircle output profile creates modified wavy output profile due to user pressing down on the flexible TPU partition moving the embedded channel. }
  \label{fig:teaser}
\end{teaserfigure}

% \received{20 February 2007}
% \received[revised]{12 March 2009}
% \received[accepted]{5 June 2009}

%%
%% This command processes the author and affiliation and title
%% information and builds the first part of the formatted document.
\maketitle

\section{Introduction}

Optical fibers
% , integral to high-tech industries such as telecommunications and medical imaging due to their ability to guide light with minimal loss, 
harbor potential in HCI applications beyond their traditional usage in telecommunications and medical imaging.
% By leveraging these fibers in desktop Fused Deposition Modeling (FDM) 3D printers, we introduce a streamlined approach to integrate sophisticated technology into everyday objects. 
% This integration not only makes cutting-edge technology accessible but also functional within the familiar confines of user-customizable objects, promising a new era of interactive devices fabricated at home or in small workshops.
% Despite the widespread use, embedding them seamlessly and cost-effectively into 3D-printed objects presents significant challenges, especially when using accessible fabrication methods like FDM. 
While previous works have explored light delivery through high-resolution  techniques for 3D printed transparent channels \cite{printed_optics_2012, clr_printed_optical_fibers_2014},
% However, these methods often fall short in terms of accessibility and cost. 
% Our research proposes a modular approach utilizing FDM, which supports uninterrupted channel printing and assembly of complex structures, 
the use of low cost, desktop Fused Deposition Modeling (FDM) 3D printing to investigate what transparent channels can afford in interactive devices has been underexplored.
We envision FDM-printed optical channels without specialized hardware, lowering access barriers for diverse users.

% offering a solution that balances functionality with user-driven customization.
%\ergun{Suggestion: While prior work in HCI has explored light delivery through high-resolution 3D printing techniques \cite{printed_optics_2012, clr_printed_optical_fibers_2014}, our approach achieves similar integration using accessible, low-cost desktop Fused Deposition Modeling (FDM) printers. This enables the embedding of optical fiber functionality—such as uninterrupted light channels—directly into complex structures, without requiring specialized equipment.}

We present a framework that applies Voronoi-based segmentation to 3D models, enabling modular fabrication of objects with continuous transparent channels embedded within each printed segment. These channels effectively route and guide light through FDM-printed structures while maintaining optical continuity (Figure~\ref{fig:teaser}(a-d)).
To develop this framework, we first examine the light transmission capabilities of transparent materials available for FDM, i.e., PolyLactic Acid (PLA), and Thermoplastic PolyUrethane (TPU), %, and PolyEthylene Terephthalate Glycol (PETG)), 
focusing on how geometric parameters influence light propagation, including channel length, diameter, and curvature. 
Additionally, we examine how printing parameters such as layer height, line width (which affect the diameter of each printed line within the channel),  and raster angle (which determines the orientation of the printed lines) affect the light transmission.
%Building on top, we develop Voronoi-based segmentation as a modular fabrication strategy for printing objects with different functional components that can be either structural and/or light transmission-capable components as shown in .


Finally, employing the modular fabrication approach, each segmented partition is separately printed and assembled 
% on standard desktop FDM printers, %\ergun{
% our modular fabrication strategy allows to fabricate objects that can 
enabling complex interaction design around light, as a controlled signal source for communication. %, through routing and channeling.
% the manufacturability and structural integrity of the final assembled object. 
% By breaking down complex geometries into manageable parts, we effectively harness the capabilities of FDM to produce intricate designs that were previously challenging to achieve with traditional 3D printing methods.
%Introducing assembly-aware segmentation of complex configurations, 
% our method enhances both the affordability and accessibility of advanced 3D printing applications. This 
Our modular fabrication strategy not only 
simplifies the fabrication process on standard desktop printers but also empowers users to experiment with different materials and interaction design techniques.
% , promoting innovation and personalization in the creation of interactive devices.
% Our FDM-based approach uniquely facilitates a wide array of existing and new interaction primitives enabled by embedded optical channels. 
% Moreover, because of this modular fabrication strategy, 
Parts of the target object can be printed in different materials for extent of light transmission and mechanical properties of the desired material, 
% for instance, by embedding channels in flexible materials like Thermoplastic Polyurethane (TPU), we can dynamically reroute these channels through end-user interactions such as 
affording different light re-routing interactions (e.g., pushing, squeezing, or rotating). % to act as I/O blocks through dynamic rerouting and direct channeling techniques. 
% connectors for multiple channel configurations. 
Complex routing or multiplexing within the object %can enable different mappings of input and output, 
can enable flexibility for smart functional design 
% the development of smart objects that can adapt their function based on user needs, 
such as adjusting light paths or altering the object's interactive capabilities in real-time (e.g., Figure \ref{fig:teaser}(e)). 
% While optical fibers have a wide range of applications, embedding them seamlessly and in cost-effective ways into 3D-printed objects remains a significant challenge, particularly when using accessible fabrication methods such as Fused Deposition Modeling (FDM). While inserting optical fibers into fabricated objects is possible, it presents significant challenges due to geometric constraints and integration complexity. 
% The difficulty lies in defining precise light paths within the object when only the endpoints are specified as inputs. Previous work \cite{printed_optics_2012,brockmeyer2013papillon, clr_printed_optical_fibers_2014} has leveraged advanced routing algorithms and high-resolution printing techniques such as PolyJet to address this issue. By routing light through light tubes within objects, Willis et al.\cite{printed_optics_2012} have achieved interactivity such as push, rotary sensing, and light relay for displays. However, FDM-based approaches are inherently limited by geometric constraints. For optical channels to be practical, they must not only enable effective light propagation but also be manufacturable within FDM’s constraints while ensuring the structural stability of the printed object. 
% % within fabricated objects can become difficult especially 
% % Within the personal fabrication domain, this passive ability of light transmission has been explored through creation of 
% % While this work proves the viability of embedding optical elements in 3D prints and a domain for interesting applications, the fabrication process itself uses PolyJet 3D printing, resources and equipment for which are expensive and hence inaccessible to the general population. 
% We propose a modular approach to achieve interactivity through optical channels using FDM that allows uninterrupted channel printing. We develop a script that enables partitioning a 3D model with complex embedded channels such that each partition can be printed separately to account for each continuous channel printed in the same plane, and then assembled to create the final object. The ability to create modular partitions not only enables fabrication on desktop FDM printers, but provides further freedom to makers to print different partitions with different materials, leveraging their properties for interactivity.
% We are hence able to print channels with complex geometries and orientations within 3D objects easily through this modularity. We further explore these complex geometries of channels, and how their geometry affects the output light compared to the input light. To improve accessibility
In sum, we contribute:

    \begin{itemize}[noitemsep, leftmargin=*]
        \item Systematic evaluation of characteristics of transparent filaments for FDM printing for their potential as optical channels,
        %\item Techniques to route and channel light to allow for complex N:M input-output mapping,
        % within an FDM printed interactive devices,
        \item Voronoi based computational geometry processing method to partition 3D objects with embedded optical channels for in-place printing, and
        % \item Accessible FDM techniques of embedded 3D printed optical channels within 3D printed objects through insertion as well as segmentation for in-place printing, that work with both single and multi-extruder printers.
        \item Design space of light routing and channeling from input to output, for  sensing, and dynamic information display. % afforded by material characterization of transparent FDM filaments. 
        % channel geometry, orientation, and modularity
        % \item Exemplary applications demonstrating the interactivity proposed in the design space 

    \end{itemize}


% \begin{figure*}[h]
%     \centering
%     \includegraphics[width=\linewidth]{figures/Workflow.png}
%     \caption{\jk{It's tedious to interpret what each block describes, each block are NOT is the same level of the idea. Redraw or just remove. }}
%      \Description{} 
%     \label{fig:workflow}
%     \vspace{-2mm}
% \end{figure*}
 
\section{Background}

% We review technologies at the intersection of developing optics devices and the fabrication of interactivity within HCI that inform our approach towards FDM printed optical devices. %. We explore prior methods to develop optical devices for sensing and dynamic display.
% , particularly in contexts such as textiles, interactive surfaces, and 3D printed interfaces. 
%We also highlight how prior research has enabled interactivity in passive 3D printed objects through optical encoding, material properties, and embedded geometries. These insights inform our approach to designing low-cost, multi-input/output systems using FDM. Finally, we examine various fabrication strategies for multi-component and modular objects, focusing on embedding optical and electronic components during or after printing. 

\subsection{Optical Devices in HCI}
Optical elements have been used extensively in Computational Fabrication for HCI research to enable sensing and information transmission. Examples include optical fiber touch sensors 
for interactive tables \cite{go2012object} and touch screens \cite{holz2013fiberio}, room-scale displays \cite{swaminathan2020optistructures}, textiles \cite{olwal2018braid, hashimoto2013lightcloth}, and grasp-sensitive surfaces \cite{wimmer2010flyeye}. To enhance flexibility, novel methods such as cellulose-based applications \cite{guridi2023light} and transparent elastomers \cite{yao2014integrating} have been explored. In the 3D printing domain, \textit{Printed Optics} \cite{printed_optics_2012} demonstrated the creation of optical channels through PolyJet printing for sensing and simple to advanced, curved displays \cite{brockmeyer2013papillon}, AR mapping \cite{fibar_2020}, and optimized routing \cite{clr_printed_optical_fibers_2014} of printed channels.
Although PolyJet printing provides significant advantages in creating arbitrarily shaped optical channels for assembly-free interactive device fabrication,
% easy component embedding, it is expensive and less accessible. In contrast, FDM is low-cost and accessible for desktop 3D printing, enabling makers and designers to design and develop their custom optical devices. 
% In this work, 
we demonstrate how FDM can be used to create optical channels for display and sensing. Furthermore, due to the fine-grained control over printing parameters and channel geometry, we show how light channeling can be varied through parameter changes. %, as shown by the material characterization.


\subsection{Optical Interactivity in Passive 3D Objects}
% Beyond creating optical pathways for light, 
% Researchers have explored methods for embedding interactivity into 3D printed objects—enabling input and output behaviors without relying on traditional electronics. 
% Prior work has incorporated photochromic \cite{jin2020photo} and thermochromic \cite{moon20243d} materials into 3D printed surfaces, allowing them to change color in response to light or heat. %These effects have also been combined in shape memory polymers to create context-aware visuals, such as QR codes that become visible only under specific conditions. Additionally, 
Interactivity is embedded through geometric or material encodings, only detectable by machines such as subsurface air cavities readable via structured light \cite{li2017aircode}, layered barcodes %using alternating material patterns 
\cite{maia2019layercode}, and infrared-fluorescent inlays for markerless tracking \cite{dogan2023brightmarker}. %Recent work has also used 
Lenticular structures have been used to produce view-dependent images %that shift displayed information based on the user’s perspective 
\cite{zeng2021lenticular}.
Surface anisotropy enables directional optical patterns (e.g., barcodes or icons) that are visible under specific lighting angles \cite{ma2023anisotag, Lu2025LumosX}.
% Internal channel design has enabled passive mechanical computation in 3D printed objects. Researchers have printed 
Embedded pipes act as light guides, or housings for post-processing elements \cite{series_of_tubes_2014}. %Air-driven logic gates \cite{airlogic_2022}, fabricated entirely from 3D printed geometries, can process user input and trigger physical or auditory outputs.
% Touch input has also been explored by embedding grid patterns beneath curved surfaces, allowing for high-resolution, multi-touch interactions directly on the object’s shell \cite{capacitive_touch_2024}.
While these prior works demonstrate approaches for embedding interactivity into 3D printed objects utilizing optics theory, most focus on isolated optical effects or tagging mechanisms. 
% Either the optical input and output are in one plane for almost parallel mapping\cite{printed_optics_2012}, or input to output mapping resolution is limited to 1:1 or 1:N \cite{series_of_tubes_2014}. In this work, we reframe optical channel routing as a design and interaction space, providing low cost ways of fabricating complex N:M input-output mapping.

In this work, we reframe optical channel routing as a design and interaction space, providing low cost ways of fabricating complex N:M input-output mappings.
Unlike previous methods where optical I/O is arranged in a single plane for nearly parallel mapping \cite{printed_optics_2012}, or limited to 1:1 or 1:N through auto-routing algorithms \cite{series_of_tubes_2014}, our approach uses the generated channels to segment the entire object into discrete and printable modules which allows users the flexibility for post-print modifications and reconfiguration. %To support FDM printing, channels must lie in-plane to create partitions.


\subsection{Fabricating Multi-Component Objects}
Printing optical channels for sensing and display within FDM 3D objects requires multi-part fabrication. 
Modular fabrication is an approach for enabling multi-component 3D printing,
% enhances design flexibility by 
decomposing a design into smaller 3D printed parts that are later assembled (e.g., \cite{series_of_tubes_2014}). 
In this way, integration of materials with various properties becomes relatively affordable,
compared to 
% technique enables designs that would be challenging  with 
conventional fabrication methods due to %. By dividing designs into smaller modules, designers can bypass 
printer size limitations and complexity constraints, making it possible to build large-scale or highly detailed objects through incremental assembly \cite{sun2021flextruss, huang20233d,moya2024single}. 
Modularity simplifies creating complex moving structures, separating functional and aesthetic sub-parts \cite{faruqi2023style2fab}, enabling interactive behaviors \cite{chen2021pneuseries}, %Additionally, modularity 
and accelerating prototyping by allowing parallel fabrication \cite{yildiz2024modular}, %when designs are divided into independent components, these parts can be produced simultaneously on multiple machines, 
significantly shortening overall build time \cite{nisser2021laserfactory, noodleprint2025}, while also producing structurally robust prints that topologically interlock the pieces \cite{voronoodle2022}. Modularity also allows multiple printers to print the object at the same time--- Voronoi-based schemes have been employed to divide and schedule printing work between different printers \cite{safezone2024, safezonestar2025}. Our work adopts modular FDM fabrication for optical devices to iteratively build complex objects, precisely place materials, and embed optical components during printing.
% , we utilize the modular design approach combined with hybrid assembly processes to ease the 



% \input{tex/2.1_Working_Principles}




\section{Working Principles}

\begin{figure}[h]
    \centering
    \includegraphics[width=0.98\linewidth]{figures/tir.png}
    \caption{Light incident in a denser medium at an angle greater than the critical angle reflects within the medium.}
     \Description{A conceptual diagram of TIR. y} 
    \label{fig:tir}
    \vspace{-2mm}
\end{figure}
%This section outlines the principles behind light transmission in FDM-printed optical channels.
% While we do not introduce new optical physics, we identify unique fabrication constraints and affordances in FDM that influence light routing and channeling. %We detail how factors like printing orientation, channel geometry, and material choice affect light behavior in transparent filaments. We also introduce the distinction between routing (path of light) and channeling (intensity control), and present our experimental setup and findings that inform how design parameters impact optical performance in modular, interactive 3D-printed systems.
\begin{figure*}[h]
    \centering
    \includegraphics[width=\linewidth]{figures/orientation.png}
    \caption{Light transmitted when the raster is parallel, while blocked when raster is perpendicular. To ensure printing channels such that the channel raster is horizontal and continuous, we need to print each channel in one plane.}
     \Description{A visual comparison showing how the orientation of the printed raster lines affects light transmission. On the left, light passes through the channel when the raster is aligned parallel to the light path. On the middle, light is blocked when the raster is perpendicular to the light path, no light is visible through the output side. On the right shows the FDM process of printing optical channels in layers on a horizontal print bed. } 
    \label{fig:printing_orientation}
    \vspace{-2mm}
\end{figure*}


\subsection{Light Transfer in Optical Fibers \& FDM Transparent Channels}
\label{sec:light_transfer}
% \jk{what are novel theories and principles that we newly identified, compared to existing optics principles and knowledge known from 3D printed optics?} \hd{in terms of physics principles, we do not have anything than what exists, I suppose the novelty is in the modular parts that can be used for interaction with the embedded channels, but I don't think that fits in this section.}
% Optical fibers consist of an inner core and an outer cladding. The core is made of a material with a higher refractive index than the cladding. Optical fibers function based on the phenomenon of \textit{total internal reflection} (see Figure
% \ref{fig:tir}). In this phenomenon, when light enters the channel at angles greater than a certain angle known as the critical angle (i.e. the angle at which light is refracted at 90\degree), the difference in the refractive indices of the two materials causes the light to bounce back into the channel. This phenomenon repeats at all core-cladding interfaces enabling light transmission. %through internal reflection.
%\hd{add graphic of total internal reflection}
Traditional optical fibers guide light via \textit{total internal reflection} (see Figure \ref{fig:tir}) at the boundary between a higher refractive index inner core and a lower-index outer cladding. 




% \hd{try to achieve the same length as printed optics - they seem to only have tested upto 100mm}
Our empirical findings indicate that light can be visibly transmitted through a filament or transparent 3D printed channel over short distances typical in desktop 3D printing. However, the angle of light incidence to the `raster' of the channel i.e. the direction of the extruder movement, affects how much light passes through the printed channel. When the raster lines are perpendicular to the incident light, reflection loses most of the light while traveling (Figure \ref{fig:printing_orientation} (middle)). However, if the raster aligns parallel with the incident light, most of the light transmits (Figure \ref{fig:printing_orientation} (left)).


To be able to transmit light through FDM printed channels, the channel needs to be printed in continuous strokes (see Figure \ref{fig:printing_orientation} (right (3))), i.e. the channel needs to lie in the XY plane. If the channel is in any other orientation, the channel is disrupted while printing layer by layer (see Figure \ref{fig:printing_orientation} right (1-2)). %due to the print strokes are disrupted due to layera channel lies vertically within a model, FDM process prints those lines layer by layer, creating interfaces every layer that reflect incident light.
This planar routing motivates our Voronoi-based segmentation method, which decomposes complex objects into planar modules. While less appealing than PolyJet, that continuously cures tiny liquid droplets blurring the layer boundaries,
% solidify layers are continuous, in that, the channels do not have layers formed due to the nature of FDM printing, and hence come close to the actual optical fiber behavior, 
the layers introduced in the channels through FDM present interesting parameters for light transfer which can enable controlled transmission and blocking for binary signaling (0 and 1) using one material.
%\hd{need an example of controlled transmission and blocking: print}
In that sense, a single transparent filament can print both the channel and the shell (Figure \ref{fig:blocking}).

\begin{figure}[h]
    \centering
    \includegraphics[width=0.8\linewidth]{figures/blocking.png}
        \vspace{-2mm}
    \caption{Object printed in same transparent material only transmits light where the raster aligns with the incident light direction and not where the raster is perpendicular. }
    \Description{Photograph of an object illuminated by an external light source. Light travels visibly through the central portion of the object where the raster direction is parallel to the light. Everywhere else, light is blocked and the output is dark except for the central portion of the object.} 
    \label{fig:blocking}
    \vspace{-4mm}
\end{figure}



% \subsection{Light Routing vs. Channeling}
% %Within fiber optics,
% \textbf{Routing} determines \textit{where} the signal travels, and \textbf{Channeling} determines \textit{how} the signal travels. %\cite{singh2008minimum}. 
% %For example, Palma et. al. \cite{capacitive_touch_2024} use capacitance as the signal and \textit{route} wires across the 3D volume. 
% In this work, we consider routing as the paths followed by the embedded optical channels between the designated input and the output sides of an object. \textit{Routing} involves control parameters such as the path of a continuous channel, as well as interactions that can connect or multiplex different channels for light transmission. 
% We consider channeling as ways in which we can control the light transmitted through the channel. \textit{Channeling} involves controlling the amount of light that will be transmitted which is based on channel geometry parameters, and light wavelength, as well as its movement while being input. We present various routing and channeling techniques for static and dynamic interactivity, forming the design space for modular fabrication of optics in FDM.  

\subsection{Material Considerations and Evaluations}
\label{sec:eval}
As one of our motivations is to lower the barrier in creating 3D printed channels to afford various interaction spaces, we consider off-the-shelf transparent materials that are commonly used in FDM printing at the time of this work: PLA, PETG, and TPU. Although PETG showed promising clarity under the lux meter, its sensitivity to temperature and cooling speed led to unpredictable and mostly milky channels. Thus, we focus on PLA and TPU. 


\begin{figure}[h]
    \centering
    \includegraphics[width=0.9\linewidth]{figures/eval_setting.png}
        \vspace{-3mm}
    \caption{ (Top and middle-left) Sample holders and measurement setup. For testing length (20mm to 100mm), thickness (1.6mm to 4.0mm), and curvature (15mm to 95mm). (Middle-right) Process of printing curved channels. (Bottom) Neopixel RGB LED and photodiode on each end.}
     \Description{Three printed channel samples designed to isolate parameters: (top left) a holder with straight channels of increasing length, (top right) a holder with identical channels of increasing channel thickness, and (middle left) a holder with constant-length channels but increasing curvature. (middle right) Slicing and sample preparation overview. (Bottom) The RGB Neopixel LEDs and photodiodes were also shown. And the setup showing the neopixel and photodiode at opposite ends of a thickness holder allow measurement of light transmission.} 
    \label{fig:eval_setting}
    \vspace{-4mm}
\end{figure}

\begin{figure*}[h]

    \centering
    \includegraphics[width=0.89\linewidth]{figures/PLA_inplace_eval.pdf}
        \vspace{-2mm}
    \caption{Light transmission through PLA in-place printed channels. Each subfigure shows the photodiode readings for red and blue light across varying parameters (a) length, (b) curvature, and (c) thickness. PLA channels show significant loss, particularly for blue light, with increasing length, tighter bend, and smaller thickness. }
     \Description{Candle stick charts showing the light transmission levels for different geometric features of in-place printed PLA channels. Each plot chart has readings for red and blue light, red light produces more response than the blue light in all readings. First plot chart shows the effect of channel length (with 2mm channel thickness, no curvature), for channel lengths of under 50mm light transmission stays mostly the same, then follows a downward trend for 50mm and longer channels. Second plot shows the effect of channel curvature radius (2mm channel thickness, 50mm length), readings increase as the curvature radius increases. Third plot shows Channel thicknesses (with 50mm channel length, no curvature), readings for thickness 1.6 mm shows a significant loss in transmission while 2mm or more than 2mm thickness values have similar light transmission values.} 
    \label{fig:pla_inplace}
    \vspace{-2mm}
\end{figure*}

\begin{figure*}[h]

    \centering
    \includegraphics[width=0.89\linewidth]{figures/TPU_inplace_eval.pdf}
        \vspace{-2mm}
    \caption{Light transmission through TPU in-place printed channels. Each subfigure shows the photodiode readings for red and blue light across varying parameters (a) length, (b) curvature, and (c) thickness. Like PLA, TPU channels show sharp declines in transmission, particularly for blue light, with increasing length, tighter bend, and smaller thickness. }
     \Description{Candle stick plots showing the light transmission levels for different geometric features of in-place printed TPU channels. Each chart has readings for red and blue light. First  plot shows the effect of channel length (with 2mm channel thickness, no curvature), red light produces similar readings for all tested channel length where blue light readings stays the same for channel lengths of under 60mm, then follows a downward trend for 60mm and longer channels. Second plot shows the effect of channel curvature radius(2mm channel thickness, 50mm length), readings show variance in lower curvature radius setting and transmits less light on average. Readings increase as the curvature radius increases. Third plot shows Channel thicknesses(with 50mm channel length, no curvature), readings with blue lights for thickness 1.6 mm shows a significant loss in transmission while 2mm or more than 2mm thickness values have similar light transmission values.  Red light readings are higher than blue light readings and mostly unaffected by the channel thickness.} 
    \label{fig:tpu_inplace}
    \vspace{-2mm}
\end{figure*}


\begin{figure}[h]

    \centering
    \includegraphics[width=0.8\linewidth]{figures/layerheight.pdf}
    \vspace{-2mm}
    \caption{Light transmission through in-place printed channels with varying layer heights and line widths, various materials. Note the Y-axis starts from 500. %\hd{replace with eval of inplace printed samples with different layer heights and line widths}
    % Each material is consistent in high photodiode readings by both red and blue lights regardless of their cross-sectional dimensions.
    % There is minimal impact of resolution on light transmission. 
    }
     \Description{Grouped candle stick chart showing photodiode readings for in place channels with varying layer heights (0.2, 0.3, 0.4mm). Each group includes red and blue bars. The y-axis starts at 500.  The red line is typically higher than the blue line. There is an upward trend from 0.2 to 0.4 in both materials.} 
    \label{fig:layerheight}
    \vspace{-4mm}
\end{figure}


% ABS is excluded due to its harmful fumes and the availability of safer and more sustainable alternatives. 
% Since we anticipated multimaterial and modular printing, 
% We selected them based on their different material properties, such as flexibility and clarity of material, as tested with a lux meter.
% While PETG showed promising clarity under the lux meter, we found that printing temperature and cooling speed had an effect on the transparency, and multi-material printing resulted in mostly milky channels. 


%Evaluations for printed and inserted PETG channels can be found in Appendix \ref{app:a}. 
% Hence, for the rest of the evaluations we only focus on PLA and TPU.

% As compared to PolyJet printing, FDM printing of transparent filaments typically results in lower clarity. However, this can be leveraged as a control parameter because 
% While transparent filaments present lower clarity compared to photopolymer for PolyJet as in prior works, FDM allows precise control over the thickness of each line printed in the channel, the degree of light transmission can be tuned to create embedded thresholds for sensing applications. For example, by changing the geometric parameters for channels within a single object, the amount of light transmitted by each channel can be controlled. % a proximity sensor could be designed so that light only reaches when the user pushes or bends the optical channel slightly.  \hd{give examples madelein: like this??? or do you want an actual example i should cite?}


% We conducted experiments by using both in-place printed channels and separately printed, inserted channels to test the effect of different geometrical parameters on light transmission. % which is $0.4$mm, provides the most amount of light transfer when compared to other layer heights, although the difference is minimal. 
% Therefore, we have maintained this configuration, and 
% We varied the length, thickness, and curvature of the channel to assess its limitation and used $0.4$mm layer height i.e. the maximum afforded by commonly available FDM extruder nozzle diameter. Lastly, because FDM printing provides fine-grained control over each individual line in a channel, we evaluated the effect of layer height and line width (i.e. the cross section of each line). 

\subsubsection{Experiment Settings.}

To understand how design affects performance, we tested several geometric parameters: channel length, curvature, thickness, and layer resolution (layer height and line width). For each parameter, we printed sample holders with in-place printed channels, as seen in Figure \ref{fig:eval_setting}. 
The length holder tests channel lengths from 20mm to 100mm in 10mm increments (thickness constant at 2mm). The thickness holder tests channel thickness (diameter) from 1.6mm to 4mm in 0.4mm increments (length constant at 50mm). The curvature holder tests the channel curve radii from 15mm to 95mm in 10mm increments (length and thickness constant at 50mm and 2mm, respectively). Finally the layer height/line width tests 0.2mm, 0.3mm, 0.4mm layer heights with 0.21mm, 0.31mm, 0.41mm line widths respectively\footnote{Bambu Studio gives an error if the line width is less than or equal to the layer height}, with overall length and thickness set to 50mm and 2mm respectively. 

We used Bambu P1S with AMS for printing with just PLA (both shell and channel) and Flashforge Creator Pro for all samples involving TPU. We used Overture transparent TPU and Giantarm transparent PLA for printing the optical channels. For the printing of all channels, we used concentric 100\% infill, an option available on most slicing softwares (we used Ultimaker Cura for Flashforge and BambuStudio). By extending the channels out of the holders, we manually cut off the edges (shown in red dashed lines in Figure \ref{fig:eval_setting} (middle-right)) with printed paths perpendicular to the direction of light. Detailed printing parameters can be found in Appendix \ref{app:params}

\subsubsection{Procedure and Apparatus.}
Each printed channel was evaluated by projecting light from a programmable RGB LED (Neopixel / WS2812B) and measuring output with a photodiode (BPW34) on the other end (Figure \ref{fig:eval_setting} (bottom)). For circuit diagram, see Appendix \ref{app:ckt}. Light intensities for red (255,0,0) and blue (0,0,255) were tested separately at maximum brightness, and photodiode readings range from 0-1023, corresponding to 0-5V. Green was ommitted because it has a broader emission spectrum that overlaps the red and blue ends of the visible light spectrum. 5 analog values for both red and blue lights were noted for each channel in the samples, measuring a total of 3 samples for each experiment. The  charts are created using the final total of 15 data points for each test variable with the whisker ends pointing to the maximum and minimum values.


\subsubsection{Results.} Following are the findings from the evaluations:
\label{sec:results}

\begin{enumerate}[noitemsep, leftmargin=*]
    \item\textbf{ As the length increased, channels exhibited a noticeable decline in transmission.}
As the length increased from 20mm to 100mm, the photodiode readings decreased sharply, especially for blue light which dropped to near zero at longer lengths (Figure \ref{fig:pla_inplace}(a) and Figure \ref{fig:tpu_inplace}(a)). However, for the red light, TPU channels did not show much difference in transmission.

\item \textbf{Channels with tighter curves or smaller bend radii showed a reduction in transmitted light.}
The drop was most significant between 15mm and 35mm radii, beyond 35mm the transmission stabilized (Figure \ref{fig:pla_inplace}(b) and Figure \ref{fig:tpu_inplace}(b)). 

\item \textbf{Channel thickness had less impact on light transmission.} 
The channel with 1.6 mm thickness transmitted much less light than the thicker ones, particularly in the blue spectrum (Figure \ref{fig:pla_inplace}(c) and Figure \ref{fig:tpu_inplace}(c)). Blue light was affected more than red, likely due to its shorter wavelength.

\item \textbf{Increasing layer height and line width resulted in increased transmission.} At 0.2mm, both PLA and TPU showed high variability and lower transmission (Figure~\ref{fig:layerheight}). By 0.4mm, both exhibited improved transmission and tighter error bars, suggesting reduced scattering due to fewer layer transitions.

\end{enumerate}


Overall, the evaluations can help us determine the geometrical and printing, as well as light wavelength parameters to control based on the desired effect.



\section{Channel Design and Object Segmentation}

% We explore various ways in which FDM can be used to fabricate interactive objects with optical channels. These techniques include hybrid ways, and printing-in-place ways of embedding channels within 3D objects.

%This section explores how FDM can be used to fabricate interactive objects embedded with optical channels. 
% We present two complementary strategies for incorporating light-guiding paths within FDM-printed parts: print-in-place method that integrates channels directly during printing, and hybrid fabrication approach that allows for the insertion of channels after printing.
% To enable precise light routing, we introduce 
%geometric modeling tools and custom algorithms that support channel creation, model segmentation, and planar alignment for efficient fabrication. Key techniques include 
% curve-based Voronoi decomposition for segmenting complex shapes, planar channel drawing interfaces, and parametric design controls for adjusting curvature, scale, and orientation. Tools are implemented using Blender’s modeling environment \cite{blender2025} and extended with Voro++ \cite{rycroft2009voro} for volumetric decomposition.


\begin{figure}[h]
    \centering
    \includegraphics[width=0.9\linewidth]{figures/channelalign.png}
        \vspace{-2mm}
    \caption{Planar (a) and complex (b) channel alignment}
    \Description{(a) simulation of 4 x 4 planar channels in a curved cuboid, (b) simulation of channels in random planes} 
    \label{fig:alignment}
    \vspace{-3mm}
\end{figure}

\begin{figure}[h]
    \centering
    \includegraphics[width=0.9\linewidth]{figures/voro.png}
        \vspace{-2mm}
    \caption{Voronoi algorithm on 2D/3D points and curves.}
    \Description{A grid illustrating the application of the Voronoi algorithm to 2D and 3D inputs, organized by input type (Points vs. Curves) and dimension (2D vs. 3D). The top row shows 2D examples: sparse points generating a standard Voronoi diagram, and curves generating curved Voronoi partitions. The bottom row shows 3D examples: a cube with scattered points, resulting in volumetric Voronoi cells, and curved 3D paths generating corresponding 3D Voronoi partitions visualized as colorful, semi-transparent regions.} 
    \label{fig:voronoioverview}
    \vspace{-2mm}
\end{figure}


\begin{figure*}
    \centering
    \begin{tikzpicture}[node distance=2.5mm]
    \node (s1) [inner sep=-0.75mm,align=center, on grid]{\includegraphics[width=0.11\linewidth]{figures/free_form_bunny/point_select.jpg}};
    \node (s2) [inner sep=-0.75mm,align=center, on grid, right=25mm of s1]{\includegraphics[width=0.11\linewidth]{figures/free_form_bunny/point_select_line.jpg}};
    \node (s3) [inner sep=-0.5mm,align=center, on grid, right=25mm of s2]{\includegraphics[width=0.11\linewidth]{figures/free_form_bunny/point_select_canvas.jpg}};
    \node (s4) [inner sep=-0.5mm,align=center, on grid, right=25mm of s3]{\includegraphics[width=0.11\linewidth]{figures/free_form_bunny/point_select_canvas_drawn.jpg}};
    \node (s5) [inner sep=-0.5mm,align=center, on grid, right=27mm of s4]{\includegraphics[width=0.11\linewidth]{figures/free_form_bunny/bunny_first_curve.jpg}};
    \node (s6) [align=center, on grid, right=25mm of s5,yshift=-3mm]{\includegraphics[width=0.11\linewidth]{figures/free_form_bunny/bunny_all_curves.jpg}};
    \node (s7) [inner sep=-0.25mm, align=center, on grid, right=25mm of s6]{\includegraphics[width=0.12\linewidth]{figures/free_form_bunny/bunny_photo.jpg}};
    \node (canvcreation) [below= of s2] {Canvas Creation};
    \node (curvedraw) [below= of s4, xshift = 16mm] {Curve Creation};
    \draw[->, thick] (s1) -- node[align=center, below] {Finding\\Axis} (s2);
    \draw[->, thick] (s2) -- node[align=center, below] {Canvas\\Rotation} (s3);
    \draw[->, thick] (s4) -- node[align=center, below] {Finalize\\Curve}  (s5);
    \draw[->, thick] (s5) edge[bend right=25] node[above] {Next Curve Creation}  (s1);
    \node (canvas) [draw,inner sep=0mm,rounded corners=1mm, fit=(s1)(s3)(canvcreation)]{};
    \node (curves) [draw,inner sep=0mm,rounded corners=1mm, fit=(s4)(s5)(curvedraw)]{};
    \draw[->, thick] (canvas) -- node[align=center, below] {Drawing\\Curve} (curves);
    \node (loop) [draw,inner sep=1mm,rounded corners=1mm, fit=(s1)(s5)(canvas)(curves)]{};
    \node (looplabel) [below= of loop] {Design Iteration};
    \node (res) [draw,inner sep=1mm,rounded corners=1mm, fit=(s6)]{};
    \node (reslabel) [below= of res,yshift=-2mm]{Resulting Channels};
    \node (res1) [draw,inner sep=1mm,rounded corners=1mm, fit=(s7)]{};
    \node (reslabel1) [align=center,below= of res1,yshift=-1mm]{3D Printed\\and Assembled};
    \end{tikzpicture}
    \vspace{-3mm}
    \caption{Overview of manual channel creation}
    \Description{A four-part figure illustrating the workflow for manually drawing and embedding optical channels using Blender. First part of the figure shows Canvas Placement: a 3D model—specifically a low-polygon bunny—with two points selected on its surface: one at the base of the bunny and another at the head. A 2D canvas is placed between these two points, having a single axis defined by the selected two points, it can rotate around this axis. Once the rotation is defined the canvas is placed. The canvas acts as a drawing plane for the light path. The second part of the figure shows Curve Drawing: a user-defined curve drawn on the 2D canvas. The curve represents the intended optical path, with visible control points and a smooth Bézier trajectory arching through the bunny’s body.The third part of the figure shows Channel Path Visualization: how the drawn curves are transformed into 3D optical channels. The curves are embedded within the bunny model’s geometry, showing translucent tubes routed internally from the base to the head. The channels are visualized using semi-transparent material, making the path clearly visible inside the solid body. The last part of the figure shows the printed bunny with optical channels placed on top of a light source, showing the various ends of the channels from the bunny lit up.}
    \label{fig:manual_channel}
    \vspace{-3mm}
\end{figure*}
%Together, these methods allow for the creation of printable parts that embed custom light paths either during or after fabrication, supporting a wide range of interactive optical designs.

% \begin{figure}[btp]
%     \centering
%     \begin{tikzpicture}[node distance=2.5mm]
%     \node (s1) [align=center, on grid]{\includegraphics[width=0.18\linewidth]{figures/voronoi/2dpoints.jpg}};
%     \node (s2) [align=center, on grid, right=16mm of s1]{\includegraphics[width=0.18\linewidth]{figures/voronoi/2dvoronoi.jpg}};
%     \node (s3) [align=center, on grid, right =18mm of s2]{\includegraphics[width=0.18\linewidth]{figures/voronoi/2dcurve.jpg}};
%     \node (s4) [align=center, on grid, right=16mm of s3]{\includegraphics[width=0.18\linewidth]{figures/voronoi/2dvoronoicurve.jpg}};
%     \node (s5) [align=center, on grid, right=16mm of s4]{\includegraphics[width=0.18\linewidth]{figures/voronoi/2dvoronoicurvemergedsolid.jpg}};
%     \node (2dpoints) [draw,inner sep=0mm,rounded corners=1mm, fit=(s1)(s2)]{};
%     \node (2dcurve) [draw,inner sep=0mm,rounded corners=1mm, fit=(s3)(s5)]{};
%     \node (d1) [align=center, on grid, below= 18mm of s1]{\includegraphics[width=0.18\linewidth]{figures/voronoi/3dpoints.jpg}};
%     \node (d2) [align=center, on grid, right=16mm of d1]{\includegraphics[width=0.18\linewidth]{figures/voronoi/3dvoronoitransparent.jpg}};
%     \node (d3) [align=center, on grid, below =18mm of s3]{\includegraphics[width=0.18\linewidth]{figures/voronoi/3dcurve.jpg}};
%     \node (d4) [align=center, on grid, right=16mm of d3]{\includegraphics[width=0.18\linewidth]{figures/voronoi/3dvoronoicurve.jpg}};
%     \node (d5) [align=center, on grid, right=16mm of d4]{\includegraphics[width=0.18\linewidth]{figures/voronoi/3dvoronoicurvesolid.jpg}};
%     \node (3dpoints) [draw,inner sep=0mm,rounded corners=1mm, fit=(d1)(d2)]{};
%     \node (3dcurve) [draw,inner sep=0mm,rounded corners=1mm, fit=(d3)(d5)]{};
%     \path  node (l1) [anchor=south west, draw,align=center,on grid,rotate=90,rounded corners=1mm, left=11mm of s1, minimum width=17mm]{2D};
%     \path  node (l3) [anchor=south west, draw,align=center,on grid,rotate=90,rounded corners=1mm, left=11mm of d1, minimum width=17mm]{3D};
%     \node (c1) at ($(s1)!0.5!(s2)$) [yshift=11mm,draw,rounded 
%     corners=1mm, minimum width=33mm]{Points};
%     \node (c2) at (s4) [yshift=11mm,draw,rounded corners=1mm, minimum width=49mm]{Curves};
%     \end{tikzpicture}
%     \vspace{-4mm}
%     \caption{Voronoi algorithm on 2D/3D points and curves.}
%     \label{fig:voronoioverview}
%     \vspace{-3mm}
% \end{figure}







%For effective light guidance through complex internal geometries, it is ideal to fabricate optical channels simultaneously with the object. However, 
FDM printing cannot support arbitrary continuous paths along the Z-axis. Discontinuous channels, when printed, suffer from significantly reduced light transmissibility due to misaligned and discontinuous channels (see Section \ref{sec:light_transfer}). To support complex FDM-compatible channel design and light routing, we propose a novel design and segmentation framework that adheres to the physical constraints of FDM. We restrict the design to planar channels, i.e. channels that lie entirely within a single plane. Such channels can be aligned with the XY plane of the print bed, allowing them to be printed seamlessly within a single layer using concentric infill. If all channels lie in the same plane, even with channels at different heights within the object, a single alignment of the object suffices (Figure \ref{fig:alignment}(a)). However, in more complex scenarios where channels must lie in different planes to connect desired points within the object, a single alignment becomes infeasible (Figure \ref{fig:alignment}(b)).

To allow multiple alignment operations to be applied to the object we need to segment the object. Segmentation is done in a way such that each channel belongs to only one segment and each segment covers different parts of the object. To ensure that each segment can be manufactured, we have utilized a Voronoi-based segmentation algorithm by Panchal et. al. \cite{panchal2023curvedspacefillingtilesusing} that divides the object based on the spatial distribution of the channels (Figure \ref{fig:voronoioverview} (right column)). 
Voronoi-based segmentation uses channels as input and produces a segmented model. Each channel is assigned its own segment, which can then be individually reoriented so its channel plane aligns with the XY plane. The design and segmentation workflow is implemented inside Blender using additions to the source code to enable Voronoi operation, and additional Python scripts to provide UI control elements with the built-in sketching tools. For detailed workflow, refer Appendix \ref{app:blender}.


\begin{figure}
    \centering
    \begin{tikzpicture}[node distance=2.5mm]
    \node (s1) [inner sep=-0.75mm,align=center, on grid]{\includegraphics[width=0.18\linewidth]{figures/free_form_bunny/bunny_all_curves1.jpg}};
    \node (s2) [inner sep=-0.75mm,align=center, on grid, right=23mm of s1]{\includegraphics[width=0.18\linewidth]{figures/free_form_bunny/bunny_all_curves_decomposition1.jpg}};
    \node (s3) [inner sep=-0.75mm,align=center, on grid, right=23mm of s2]{\includegraphics[width=0.18\linewidth]{figures/free_form_bunny/bunny_all_curves_decomposition_channel_boolean1.jpg}};
    \node (s4) [inner sep=-0.75mm,align=center, on grid, right=23mm of s3]{\includegraphics[width=0.18\linewidth]{figures/free_form_bunny/bunny_decomposed_layout_curves1.jpg}};
    \draw[->, thick] (s1) -- node[align=center, below] {Voronoi\\Partition} (s2);
    \draw[->, thick] (s2) -- node[align=center, below] {Channel\\Boolean} (s3);
    \draw[->, thick] (s3) -- node[align=center, below] {Layout} (s4);
    \node (res) [draw,inner sep=1mm,rounded corners=1mm, fit=(s1)(s4)]{};
    \end{tikzpicture}
    \vspace{-5mm}
    \caption{Channels are used for partitioning the shape. Then each partition is aligned according to its ``seed'' curve such that the light channels will lie in XY plane.}
    \Description{Shows the process of how optical channels are used to partition a 3D model into printable modules. On the left side shows a Stanford bunny with 4 optical channels. An arrow labelled Voronoi Partition points to the Stanford bunny with the channels and the regions defined around the channels. An arrow labelled Channel Boolean points to the Stanford bunny where the space for the optical channels has been removed using a boolean operation to preserve the paths. Finally, an arrow labelled layout shows the 4 partitions where each partition is based on its corresponding channel seed so that the channels lie flat in the XY plane.}
    \label{fig:object_partitioning}
    \vspace{-4mm}
\end{figure}

\subsection{Channel Creation} 
Our design framework enables users to manually define channel curves through the volume of the object, as illustrated in Figure~\ref{fig:manual_channel}. The process is directly implemented in Blender, using Blender's Python API and built-in 3D sketching tools. To begin, the user first selects two endpoints on the object surface. These points define an axis, which serves as the rotation center for a temporary 2D canvas. This plane is essential for two reasons: (1) it ensures that the light path can be sketched on a single plane, making the channel geometry printable via FDM, and (2) it allows the user to route the channel through the solid parts of the object, thus avoiding holes and empty regions. User then adjusts the rotation of the canvas about this axis, allowing fine control of the channel’s orientation. 

Next, the tool asks users to sketch curves using the built-in 3D sketching tools in Blender, which allow using a plane for projecting the sketches. These sketches define the desired light path. After drawing, the canvas is removed, and the drawn path is approximated using a series of Bézier curves, which offer smooth interpolation and control over curvature. This approach allows fine-tuned routing---for instance, subtly curving a channel from the bottom of the object to reach the eyes of a bunny model (Figure~\ref{fig:manual_channel}). The curve retains the orientation information from the previously defined canvas, and it is used for orienting the segments later.

To ensure robust printability, channels are extended slightly beyond the object’s shell. This allows the extruded filament paths to terminate cleanly through manual cutting, minimizing scattering loss, as concentric infill at these endpoints naturally orients perpendicular to the incident light.

% \jk{add pseducode of the algorithm}
% \begin{algorithmic}
%     \caption{Algorithm for Channel Creation}
    
% \end{algorithmic}

\subsection{Partition Creation} 
%\hd{need to write about feedback when partitioning is not required}

To segment the object based on the designed channels, we first analyze the channel curves to determine whether partitioning is necessary. If all user-defined curves lie within a single plane, segmentation is not required. In this case, the object can simply be oriented with respect to that plane to facilitate fabrication. This planar scenario is identified and highlighted within the tool.

If the channels cannot be contained in a single plane, we first sample points along each Bézier curve by arc length to ensure a consistent resolution. These points serve as seeds for a 3D Voronoi decomposition. Voronoi is used for segmenting the object because it generates convex cells which are then used to intersect with the object boundary. Assuming the object is printable, boolean operation with a convex domain does not interfere with the manufacturability of the object. Traditional Voronoi diagrams are based on discrete points (Figure \ref{fig:voronoioverview} (left column)); however, we follow the curve-based Voronoi approach introduced by Panchal et al.~\cite{panchal2023curvedspacefillingtilesusing}, where sampled points from the same curve are grouped, and their corresponding Voronoi cells are fused to form a region representing the entire curve (Figure \ref{fig:voronoioverview} (right column)).

This segmentation process is implemented using Blender, where channel curves are modeled using its default curve tools. For volumetric decomposition, we integrate the Voro++ library~\cite{rycroft2009voro} into Blender’s Geometry Nodes system by implementing a C++ module that executes within the procedural workflow. Voro++ operates within a bounding box defined slightly larger than the input mesh to ensure that resulting partitions fully encapsulate the geometry and allow for robust Boolean operations.

After computing the Voronoi diagram, we finalize the segmentation using Boolean intersections between the object and each partition. This results in a set of mesh segments. To integrate the channels themselves, we perform an additional Boolean subtraction using swept volumes generated from the Bézier curves. This ensures that each channel is embedded cleanly within its corresponding mesh segment (Figure~\ref{fig:object_partitioning}). To orient each segmented piece, we leverage the previously generated canvases to determine the alignment of the associated curves. We then compute the rotation that aligns each canvas with the XY plane and apply this transformation to the corresponding segment.
The final design shown in Figure~\ref{fig:manual_channel} includes eight optical channels routed across the back of the bunny, along with two additional channels that route light from the base to the eyes, demonstrating the capability of the system to handle complex, multi-directional channel placement in a manufacturable manner. Algorithm \ref{algo_1} shows our process pseudocode.

% \jk{add pseducode of the algorithm}

\begin{algorithm}
\caption{Channel-Based Object Segmentation}
\label{algo_1}
\begin{algorithmic}[1]
\Require Set of Bézier curves $C = \{c_1, c_2, \dots, c_n\}$ on object mesh $M$
\Ensure Segmented object with embedded channels

\If{All curves lie in a single plane}
    \State Orient object such that curves lie in XY plane
    \State \Return $M$ without segmentation
\EndIf

\ForAll{$c_i \in C$}
    \State Sample points $P_i = \{p_1, p_2, \dots\}$ along $c_i$
\EndFor
\State Initialize 3D bounding box around $M$
\State Compute curve-based Voronoi diagram using grouped point sets \cite{panchal2023curvedspacefillingtilesusing}
\State Fuse cells from same curve group $P_i$ into region $R_i$

\ForAll{region $R_i$}
    \State $S_i \gets \text{BooleanIntersect}(M, R_i)$ \Comment{Segment mesh with Voronoi cell}
    \State $V_i \gets \text{SweepVolume}(c_i)$ \Comment{Generate channel volume}
    \State $S_i \gets \text{BooleanSubtract}(S_i, V_i)$ \Comment{Embed channel into segment}
    \State $S_i \gets \text{OrientSegment}(S_i, c_i)$ \Comment{Orient each segment with the associated curve $c_i$}
\EndFor

\State \Return $\{S_1, S_2, \dots, S_n\}$
\end{algorithmic}
\end{algorithm}

\begin{figure}[h]
    \centering
    \begin{tikzpicture}
        \node (s1) [inner sep=-0.75mm,align=center, on grid]{\includegraphics[width=0.45\linewidth]{figures/blender_interfaces/transform.jpg}\\Transform};
        \node (s2) [inner sep=-0.75mm,align=center, on grid, right= 45mm of s1, yshift=7mm]{\includegraphics[width=0.45\linewidth]{figures/blender_interfaces/one2many_params.jpg}\\1:Many};
        \node (s3) [inner sep=-0.75mm,align=center, on grid, below= 20mm of s2]{\includegraphics[width=0.45\linewidth]{figures/blender_interfaces/shape2shape_params.jpg}\\Shape:Shape};
    \end{tikzpicture}
    \caption{Custom control parameters extracted in Blender.}
    \Description{Blender interface showing design control blocks for transforming channel geometries. The Transform shows the value of Block Width: 0.1 m, Block Height: 0.1 m, Num Channels: 3, Z Axis Rotation: 0{\degree}, Scale: 0.400, 0.400, 1.000, Shear: 0.03 m, Curvature: 0, Translation: 0,0,0. The 1:Many block shows the values Cube Length : 0.5 m, Middle Point X, Y, Z: 0 m, Channel Gap: 0.5 m, The Shape:Shape block shows Channel Amount: 8 and and option to add 2 shapes.}
    \label{fig:modifier_params}
    \vspace{-2mm}
\end{figure}



\begin{figure}[h]
    \centering
    \begin{tikzpicture}[node distance=2.5mm]
    \node (a1) [align=center, on grid]{\includegraphics[width=0.18\linewidth]{figures/transforms/one2one_input.jpg}};
    \node (a2) [align=center, on grid, right=18mm of a1]{\includegraphics[width=0.18\linewidth]{figures/transforms/one2one.jpg}};
    \node (a) [draw,inner sep=1mm,rounded corners=1mm, fit=(a1)(a2)]{};
    \node at ($(a1)!0.5!(a2)$) [yshift=-8mm]{1:1};
    \node (b2) [align=center, on grid, below=20mm of a1]{\includegraphics[width=0.18\linewidth]{figures/one2many/cube_init_curves.jpg}};
    \node (b3) [align=center, on grid, right=18mm of b2]{\includegraphics[width=0.18\linewidth]{figures/one2many/cube_init_curves1.jpg}};
    \node (b4) [align=center, on grid, right=20mm of b3]{\includegraphics[width=0.18\linewidth]{figures/one2many/cube_init_curves2.jpg}};
    \node (b5) [align=center, on grid, right=18mm of b4]{\includegraphics[width=0.18\linewidth]{figures/one2many/cube_partitions.jpg}};
    \node (b) [draw,inner sep=1mm,rounded corners=1mm, fit=(b2)(b5)]{};
    \node at ($(b3)!0.5!(b4)$) [yshift=-8mm]{1:Many};
    \node (c3) [align=center, on grid, right=20mm of a2]{\includegraphics[width=0.18\linewidth]{figures/shape2shape/input_interpolated.jpg}};
    \node (c4) [align=center, on grid, right=17.5mm of c3]{\includegraphics[width=0.18\linewidth]{figures/shape2shape/partitioned.jpg}};
    \node (c) [draw,inner sep=1mm,rounded corners=1mm, fit=(c3)(c4)]{};
    \node at ($(c3)!0.5!(c4)$) [yshift=-8mm]{Shape:Shape};
    \node (d1) [align=center, on grid, below= 20mm of b2]{\includegraphics[width=0.18\linewidth]{figures/transforms/rotation_input.jpg}};
    \node (d2) [align=center, on grid, right= 18mm of d1]{\includegraphics[width=0.18\linewidth]{figures/transforms/rotation.jpg}};
    \node (d) [draw,inner sep=1mm,rounded corners=1mm, fit=(d1)(d2)]{};
    \node at ($(d1)!0.5!(d2)$) [yshift=-8mm]{Rotation};
    \node (e1) [align=center, on grid, right= 20mm of d2]{\includegraphics[width=0.18\linewidth]{figures/transforms/scaling_input.jpg}};
    \node (e2) [align=center, on grid, right= 18mm of e1]{\includegraphics[width=0.18\linewidth]{figures/transforms/scaling.jpg}};
    \node (e) [draw,inner sep=1mm,rounded corners=1mm, fit=(e1)(e2)]{};
    \node at ($(e1)!0.5!(e2)$) [yshift=-8mm]{Scaling};
    \node (f1) [align=center, on grid, below= 20mm of d1]{\includegraphics[width=0.18\linewidth]{figures/transforms/shearing_input.jpg}};
    \node (f2) [align=center, on grid, right= 18mm of f1]{\includegraphics[width=0.18\linewidth]{figures/transforms/shearing.jpg}};
    \node (f) [draw,inner sep=1mm,rounded corners=1mm, fit=(f1)(f2)]{};
    \node at ($(f1)!0.5!(f2)$) [yshift=-8mm]{Shearing};
    \node (g1) [align=center, on grid, right= 20mm of f2]{\includegraphics[width=0.18\linewidth]{figures/transforms/tranlation_input.jpg}};
    \node (g2) [align=center, on grid, right= 18mm of g1]{\includegraphics[width=0.18\linewidth]{figures/transforms/tranlation.jpg}};
    \node (g) [draw,inner sep=1mm,rounded corners=1mm, fit=(g1)(g2)]{};
    \node at ($(g1)!0.5!(g2)$) [yshift=-8mm]{Translation};
    \end{tikzpicture}
    \caption{Parameters for design primitives and transforms.}
    \Description{Diagram of design parameters. The 1:1 block shows 9 channels evenly distributed on the top to the bottom of the cube in the first part of the 1:1 block, the next part of the 1:1 block shows how the cube is partitioned with the channels. The Shape:Shape block shows a 6 input channels arranged in a line at the bottom of the cube and 6 output channels arranged in a circle on top of the cube. Similar to the 1:1 block, the first part of the image shows the location of the channels in the cube while the second part of the image shows the channels and how they are partitioned in the cube. The 1:Many block shows 4 cubes, 3 of the cubes show the cube with the channels in different orientations and the last cube shows the cube with the partitions and channels. There are 5 channels, one on each side of the cube that all converge to the last side of the cube. The rotation, scaling, shearing, and translation effects are also shown with 2 images. The rotation INSERT ROTATION HERE. The scaling shows how 9 channels are evenly distributed on the bottom of the cube and all converge to a small grid on top of the cube. INSERT SHEARING AND TRANSLATION.}
    \label{fig:blocks}
\end{figure}

\begin{figure}[h]
    \centering
    \begin{tikzpicture}
        \node (s1) [inner sep=-0.75mm,align=center, on grid]{\includegraphics[width=0.2\linewidth]{figures/printedblocks/o2m1.jpg}};
        \node (s2) [inner sep=-0.75mm,align=center, on grid, right= 20mm of s1]{\includegraphics[width=0.2\linewidth]{figures/printedblocks/o2m2.jpg}};
        \node (s3) [inner sep=-0.75mm,align=center, on grid, right= 20mm of s2]{\includegraphics[width=0.2\linewidth]{figures/printedblocks/s2s1.jpg}};
        \node (s4) [inner sep=-0.75mm,align=center, on grid, right= 20mm of s3]{\includegraphics[width=0.2\linewidth]{figures/printedblocks/s2s2.jpg}};
    \end{tikzpicture}
    \caption{Printed examples of 1 : Many and Shape : Shape routing blocks.}
    \Description{The first part of the image shows 1 channel channel of a 1:Many block illuminated. The second part of the the image shows 3 channels on different sides of the cube of a 1:Many block illuminated. The third part of the image shows 2 out of the 7 channels arranged in a line of a Shape:Shape block illuminated. While the last part of the image shows 2 out of the 7 channels arranged in a circle of a Shape:Shape block illuminated.}
    \label{fig:printedblocks}
    \vspace{-5mm}
\end{figure}

\subsection{Parametric Design}
To ease the creation of different I/O routing, we extracted parameters into a set of modifier modules, offering users control over parameters defined in the control panel in addition to manual modifications. We have implemented 3 main blocks (see Figure \ref{fig:modifier_params}) for parametric design:% Each block has different parameter space:

\begin{itemize}[noitemsep, leftmargin=*]
\item \textbf{ Transform blocks.} User can select the number of channels distributed in a square grid and change the domain size of the block. Transform blocks combine multiple transformations: Rotation, Scaling, Shearing, and Translation. These blocks can be rotated in space  for manual alignment. The rotation within the block refers to  how channel entries are rotated around the Z-axis. Scaling operation scales the endpoints of the curves. Shearing operation moves the end points of the channels on X-axis depending on the Y-position of the end points. Translation operation translates the end points in a given vector direction.
\item \textbf{1 : Many blocks.} This block enables a channel to be split into multiple branches where the end points are located on the center of the faces of a cube. The parameters include the size of the cube, middle control point of the curves for determining curvature of the channels, and the channel gap size between different curves.
\item \textbf{Shape : Shape block.} This is a modifier to interpolate new channels between two different input shapes, using the two shapes and the number of channels to be interpolated between them.% two shapes.
\end{itemize}

%These design parameters allow for expressive channel geometry, 
Figure \ref{fig:blocks} highlights the parameters and effects on channel geometry. 
The script then generates a bounding box for the channels which can be modified using external cut and extrude options. Then the Voronoi-based partitioning modifier creates the partitions of the model. Lastly, another modifier is applied to align all the partitions on a plane such that all channels are printed in the XY plane. These models can then be exported and sliced in any slicer for printing, and assembled with an adhesive as shown in Figure ~\ref{fig:printedblocks}.








% \subsection{Hybrid Fabrication for Embedding}
% Hybrid fabrication combines in-place printing of complex routes with casings that have built-in spaces for channels. Flexible materials like TPU make manual insertion of channels easier.
% % In addition to computational geometry segments for printing complex routes in-place, hybrid fabrication involves printing the casing with spaces for channels in the 3D space, and inserting the channels within the object, easier with TPU by its flexible nature.
% Channels can be printed based on the required length and thickness, then manually inserted into the object afterward (Figure \ref{fig:hybridinsertion}). In this case, the filaments themselves can also be utilized as channels, removing the need for special printing of channels. This technique leverages the modularity of already created partitions enabling parallel printing on different machines, especially when multi-extruder printers are not available. However, this technique requires careful planning of pathways as the channels cannot be inserted through extremely convoluted paths, requiring in-place printing.

% \begin{figure}[h]
%     \centering
%     \includegraphics[width=\linewidth]{figures/hybridinsertion.png}
%     \caption{By creating internal channel geometries, transparent flexible filaments (left), or channels printed with them (right) can be directly inserted}
%      \Description{Hybrid fabrication illustration showing insertion of transparent filament and a printed channel. The left part of the image shows a holder that has a curved path and a flexible filament is inserted to act as a channel. The right part of the image shows a straight printed channel being inserted to a holder with a straight path. } 
%     \label{fig:hybridinsertion}
%     \vspace{-2mm}
% \end{figure}



\section{Techniques to Afford Interaction with Optics}


%Application scenarios of this approach is wide.
% for example, one-to-one mapping of channels can be utilized as a routing display, where the input surface of the object can be placed on a conventional dynamic input display such as LCD screens, and the display surface of the 3D printed optical device would be directly mapped to the output surface. 
% A scenario of an end user playing an active interaction is possible, navigating through a dynamic motion path, linear to rotational, and creating shadows, etc. 
We introduce various interaction affordances and their proofs of concepts,
% , the type and motion of input light, material properties, and human interventions. We 
starting from a simple mapping and adding on complexities to explain the interplay of control parameters. Note that objects with channels lying in the same plane are not segmented, but printed directly (case in Figure \ref{fig:alignment}(a)). All other objects have been segmented and modularly printed.

\subsection{1:1 Parallel Mapping}
We can utilize the orientation of the channel geometry in 3D space as embedded in an object to route light from one point to another (a one-to-one mapped display routing), as shown in a toy eye reflecting changing emotions from LCD projection input \cite{printed_optics_2012}.
% shows the use of light pipes to display the location of chess pieces, as well as light up the eyes of a printed toy using a mobile projection as the source. 
Similarly, we consider each channel end to be one display point (i.e., one pixel), with multiple channel ends forming a display (see Figure \ref{fig:lighteffect}), which can be triggered by any light source.  % as the printed channels tend to transmit various signal types and luminance.
%The goal is to transmit light signals without altering its spatial structure or orientation.
%1:1 mapping technique achieves direct transmission of light source(Figure \ref{fig:lighteffect}). % to shadows when light is blocked (Figure \ref{fig:lighteffect}b) to ambient colors (Figure \ref{fig:lighteffect}c).


% \begin{figure}[h]
%     \centering
%     \includegraphics[width=\linewidth]{figures/simpledisplay.png}
%     \caption{(a,b) Sample with one-to-one channel mapping displays colors directly projected through input (c) \hd{remove}}%\jk{this figure is too large given what it tells beyond known community knowledge i.e., 3D printed optics paper}}
%      \Description{An object with embedded optical channels aligned in a grid. (a) shows a rendered image that is stopped midway through to show the path of the channels. (b) shows a printed object with boxes to show the input and output side. (c) shows the input side facing downward onto a light source displaying discrete blocks of color. The output face shows corresponding colored dots at matched positions to verify 1:1 mapping between input and output.} 
%     \label{fig:simpledisplay}
%     \vspace{-2mm}
% \end{figure}


\begin{figure}[h]
    \centering
    \includegraphics[width=\linewidth]{figures/onetoone.png}
    \caption{One-to-one light mapping without segmentation} 
    \Description{In-place printed channels without segmentation with 1:1 light mapping, with left image showing a sliced model, and right showing the actual print transmitting colored light.} 
    \label{fig:lighteffect}
    \vspace{-3mm}
\end{figure}


\subsection{Non-parallel Mapping}
Adding a layer of complexity, we introduce offsets between the input space and output planes by routing channels along paths. This enables transmitting signals to select endpoints based on the routing such as one-to-one, one-to-many, many-to-many connections, etc (e.g., Figure \ref{fig:printedblocks}). 

\paragraph{Slider Mechanism.}
A slider mechanism presents  an interaction where the channels are routed such that  a linear input is optically transformed into a semicircular output.
 % to move the input light that can be channeled through the embedded optical channels, and 
The physical translation motion of the slider is converted to a visual rotational motion.
Figure \ref{fig:slider} shows how moving the slider over static lights presents a semicircle motion of colors, with colors moving to the left, with linear sliding motion to the right.%each new color being revealed with further linear sliding motion.
%For example, Figures \ref{fig:muxing_1}a and \ref{fig:muxing_2}a show how direct routing affects what information is transmitted on the output end and in what order. In Figure \ref{fig:muxing_1}, the channels are created such that an input end point of a channel moves one space next to what it originally would have been according to Figure \ref{fig:simpledisplay}. Similarly, in Figure \ref{fig:muxing_2}, the input end points are arranged in a line while the output end points are arranged in a semicircle. Such routing results in a modified output display than what is input. % \jk{this paragraph lacks detail. fully describe what's possible in text so the reader can understand this per see not by seeing the figures or design space.}



\begin{figure}[h]
    \centering
    \includegraphics[width=\linewidth]{figures/slider.png}
    \caption{Switch slides linearly over a light strip showing a circular output light movement.}
    \Description{A slider body is placed on top of a horizontal strip of LED lights. As the slider is translated to the right along the light source, output light emerges in a rotational motion pattern from red to green to blue in a semicircular shape. This shows the conversion of linear motion of the input into a dynamic and rotational visual output. The figure illustrates how the slider's internal channels reveal three different colored light paths depending on the input.} 
    \label{fig:slider}
    \vspace{-2mm}
\end{figure}

\begin{figure}[h]
    \centering
    \includegraphics[width=\columnwidth]{figures/orientation_sense_1col.png}
    \caption{Many:1 channel creation with each channel having different length and curvature, affects amount of light transmitted enabling orientation sensing of the object.}
     \Description{The process of creating a cube demonstrating object orientation sensing using Many:1 channel creation. The first part of the photo shows a rendered cube with channels with different lengths and curvatures. The next part of photo shows an assembled version of the cube with 3 optical channels. The next part of the photo shows the exported parts and the sensing setup that included a neopixel grid and the 3 sides of the cube labeled. The next parts of the photo show a setup where each side is detected facing downwards based on the value of light transmitted.} 
    \label{fig:orient}
    \vspace{-3mm}
\end{figure}









\paragraph{Object Orientation Sensing.}
Using a many-to-one channel mapping, it is possible to detect which part of an object is facing a light source. By adjusting the geometric parameters---such as the length and curvature---of multiple PLA channels, we can control how much light each one transmits. When these channels converge at a single sensor, the amount of light received from each channel provides information about the object's orientation relative to the light source. Figure \ref{fig:orient} shows how three sides of a cube can be detected to be facing downwards (toward the light) by thresholding the light transmitted by each channel. Each channel is designed with a unique length and curvature to change the amount of light transmitted (refer Subsection \ref{sec:results} (\#1,2)). Similarly, blocking of light, and  wavelength can also be sensed through the evaluation data and appropriate thresholding.



\subsection{Motion Steering}
% We look at the motions in which the light movement become one type of input, and the motions in which it is output. 
Steering of light motion reflects a logic or adaptation that governs how output elements move in relation to input light in motion.
%For example, in Figure \ref{fig:simpledisplay}b, moving input light from left to right on one side would map to output light moving from right to left.% on the output side.
% based on the geometry of channels within the sample and movement of light. \jk{this paragraph is also not self-descriptive}

\begin{figure}[h!]
    \centering
    \includegraphics[width=\linewidth]{figures/motionsteer.png}
    \caption{Steering light in a clockwise motion gives output light movement in anti-clockwise motion.}%\jk{This figure may benefit from drawings instead. to clearly showcase input and output mapping in abstract figure.}}
    \Description{A printed object demonstrating non-parallel optical routing between input and output faces. (a) shows a diagram with the input as yellow, blue, red, and green in a clockwise direction and the output as red, blue, yellow, and green in the clockwise direction; (b) shows a light beam moved clockwise across the input face, the output beam shifts anticlockwise.} 
    \label{fig:muxing_1}
    \vspace{-2mm}
\end{figure}

% \begin{figure}[h!]
%     \centering
%     \includegraphics[width=0.9\linewidth]{figures/muxing_2.png}
%     \caption{Static routing with a line to semicircle mapping of input and output channel ends}
%     \Description{An object with input channel arranged in a line along one face and output channel ends arranged in a semicircle on the opposite face. (a) shows a diagram, a rendered image, and a printed object with a light source connected at the end. (b) shows how the routing transforms a linear motion at the input into a curved path at the output. As a light beam travels left to right along the linear input, it sequentially lights up the channels across the arc from right to left.} 
%     \label{fig:muxing_2}
%     \vspace{-2mm}
% \end{figure}

Due to the way channels are routed in ~Figure \ref{fig:muxing_1}(b), when input light is moved clockwise, the output light moves anticlockwise. %(instead of right to left as would be the case for Figure \ref{fig:simpledisplay}). 
This can also be seen from the modified color mapping from input to output (Figure \ref{fig:muxing_1}(a)). %When changed the paths of the channels, such as twisting the channel curve 90\degree, or inversing the arrangement of input channel that connects to the ends as output, linear incident light horizontally moving can present vertical movement, as seen in . 
%As seen in Figure \ref{fig:muxing_1}a, the color mapping of each color is one space next to what it would have been if the channels were not rotated 90\degree during routing. \jk{this is hard to parse. I don't get spatial understanding of this easily.} 
%Similarly, arranging input ends of the channels in a line, and output ends in a semicircle gives a mapping from a line to a curve (Figure \ref{fig:muxing_2}). If we add movement of input light along the line from left to right, we see the output light move in a semicircle from right to left.

\begin{figure}[h]
    \centering
    \includegraphics[width=\linewidth]{figures/wheelalignment.png}
    \caption{A semicircle to curve mapping enables detecting deviations from the center.}
     \Description{A driving assistance indicator prototype is shown with 7 optical channels lined up on top. When the steering wheel is centered, the output shows a red light in the middle (4th). Turning the wheel slightly shows the light move to the third channel for slightly left and the fifth channel for slightly right. Turning the wheel further to the left lit up the second optical channel.} 
    \label{fig:wheel}
    \vspace{-2mm}
\end{figure}


\paragraph{Driving Assistance Indicator.}
Using a mechanism to steer the light rotationally, we designed a lane alignment driving assistance indicator (see Figure~\ref{fig:wheel}) that translates a steering wheel's left and right movements into a display on the dashboard. When the indicator remains centered, it signifies that the vehicle is properly aligned within its lane. Significant or frequent deviations from the center suggest that the driver is making excessive steering corrections, indicating a need for steadier control.


\subsection{Multiplexing and Decoupling}
% by Dynamic Re-routing}
By connecting  channels printed in different objects,  we can create new light paths.
%Discontinuity of the channel with which an inserted optical channel can refract light paths.
% We explore how the channels that have been printed or inserted can be moved to create new channeling connections. 
These include interactions afforded through the geometry of the entire object. % and through material characteristics.
% As seen through existing work \cite{printed_optics_2012}, 
We show rotary switches to validate this hypothesis, with an internal channel printed in place (Figure \ref{fig:rotary}(a, c) that can connect different input and output channels \cite{printed_optics_2012}--- the output channels here are printed separately and inserted for testing purposes.
% . In such applications, modularity can play an important role in creating 
Plug-and-play rotary switches can be designed to afford different routing behaviors, such that one switch may have a straight channel connecting an input and output in a straight line (see Figure \ref{fig:rotary}(a, b)), while another transmits  input light at 90\degree (see Figure \ref{fig:rotary}(c, d)). While not segmented, this validates the use of FDM printed channels for plug and play interactivity. With our segmentation technique, there is potential to design non-planar switches as well.
% By rotating the switch, different connections can be made.
% We can connect different inputs and outputs of channels for a plug-and-play type effect similar to Legos.



\begin{figure}[h]
    \centering
    \includegraphics[width=\linewidth]{figures/rotary.png}
    \caption{Rotating the replaceable central switch to connect different channels for light routing.}
    \Description{A modular rotary switch mechanism consisting of a replaceable central insert and a printed base with embedded channels. (a, b) show one insert design where an internal channel routes light directly across the switch in a straight line. (a) shows just the knob with an optical channel and (b) shows the knob inserted into a base with 8 different channels with the light showing up on the opposite end of the light source. (c, d) show a different insert that redirects light at a 90° angle from the input to a perpendicular output. Similar to (a) and (b), (c) shows the insert knob while (d) shows the know inserted into the base lighting up a channel 90\degree from the input light source.} 
    \label{fig:rotary}
    \vspace{-3mm}
\end{figure}






The properties of the material can be used to encase rigid channels in a sparse space, granting some motion space for the channel.
% , similarly observed in a push sensor in prior work \cite{printed_optics_2012}.
% can also be utilized for rerouting the channels, 
To test this hypothesis, we inserted a printed PLA channel in a flexible TPU casing with low infill. We can push the casing to move the embedded channel to point to a different output, creating a push button's On/Off binary signal to be captured at the receptor (Figure \ref{fig:push}). 
Flexibility of the 3D printed materials affords a wide array of tangible interactions presented in prior fabrication research, such as shear rotation to lock and unlock a door latch~ \cite{ion2016metamaterial} , or bike handles and saddles for comfortable grip and sitting~\cite{kim2021softness}. %, where we anticipate the possibility of regenerating them while embedding sensing capabilities, to sense human interactions with physical objects using this technique.
% As we propose a modular design, including such interactions in various parts of a single object can be possible.

\begin{figure}[h]
    \centering
    \includegraphics[width=\linewidth]{figures/pushbutton.png}
    \caption{Channel embedded in flexible TPU can be pushed to redirect the light towards a sensor, acting as a push button}
    \Description{A first part of the photo shows a tactile push button fabricated in TPU with an embedded optical channel connected to a light source on one end, a photodiode could be seen on the other end. The middle part of the image shows the user pushing down on the same setup, causing the optical channel to move down. The last part of the image shows a graph where there are peaks when the push button is being pressed down and the width of each peak coresponds to the duration of the press. } 
    \label{fig:push}
    \vspace{-3mm}
\end{figure}

% While these geometries, channel orientations, and light movement provide new space of interaction, fabricating these channels in 3D space using FDM is not possible. We can create the same examples by fabricating separate channels, in TPU or PLA and inserting them, however to be able to print the channels in place, we need to develop a modular technique of fabrication where each channel is printed in the XY plane as required for light transfer. We will hence describe the different techniques of fabricating such objects with FDM, and present our modular splitting method for printing channels in place using multi-material printing.

\begin{figure}[h]
    \centering
    \includegraphics[width=\linewidth]{figures/bikehandle.png}
    \caption{Squeezing TPU printed body moves the embedded channels closer together, displaying them through openings in the cap. }
     \Description{A four part photo depicting a bike handle application using a TPU printed body and a PLA handle and cap. (a) shows the default, unsqueezed state of the TPU part of the handle with optical channels spaced apart. (b) shows the TPU part of the handle being squeezed, mechanically deforming the TPU and bringing the optical channels closer. (c) shows the entire setup with both the TPU and PLA parts unsqueezed with light going through the handle. (d) shows the TPU squeezed which makes transmits the light through the holes in the cap.  } 
    \label{fig:bikehandle}
    \vspace{-2mm}
\end{figure}


\paragraph{Squeeze to Signal Bike Handle.}
% By leveraging the modularity of 3D-printed parts, we can fabricate components—such as a bike handle holder—in flexible TPU material. 
With tested hypothesis, we moved onto in-place printing of the bike handle, with the casing in TPU, allowing the in-place printed PLA channels to move closer together or further apart (Figure \ref{fig:bikehandle}(a, b)). The resulting mechanical interaction can then be transmitted through internal channels, making the signal visible to others (Figure \ref{fig:bikehandle}(c, d)).
In this setup, different types of interactions can be mapped to distinct meanings. For instance, squeezing the handle could trigger a red light on the handlebar to signal caution. Alternatively, the design could incorporate directional input: a top press might shift a movable channel to connect with a static green-lit channel, signaling that it’s safe for others to overtake. While the channels in this example are in the same plane making segmentation  unnecessary, we can utilize the same principle for more complex routing similar to Figure \ref{fig:teaser}. %Conversely, a bottom press could engage a red-lit channel to indicate that overtaking is unsafe --- such as when the rider is preparing to turn.

\begin{figure}[ht]
    \centering
    \includegraphics[width=\columnwidth]{figures/designspace_1col.png}
    \caption[]{Design Space of Light Routing, Channeling, and User Interactions. Icons from Flaticon\footnotemark.}
    \Description{The full design space of optical interactions enabled by embedded channels. The layout includes icons and labels for static I/O routing which is further subdivided into Basic Relay (e.g., 1:1, 1:Many, Many:1) and Transforms and Non-parallel Switching (Scale, Translate, Rotate, Shape to Shape, Shear), dynamic re-routing which is further subdivided to multiplexing (e.g.,plug-and-play, building on each block, slider, rotary switch), decoupling behaviors (e.g., push buttons, touch-based occlusion and bistable switch), and lighting control (e.g., lighting motion steering, color wavelength, and light blocking)}
    \label{fig:dspace}
    \vspace{-4mm}
\end{figure}

\footnotetext{Unplug by juicy\_fish; Construction by Darius Dan; Slide, Dial, Panic button by Freepik; Slider by manshagraphics; Find by howcolour, Ocean by IconBaandar.}

\subsection{Design Space}

Using this interplay of light, routes, and user interactions, we map out the design space for FDM optical pathways (Figure~\ref{fig:dspace}).
% \subsubsection{Static Interactivity}
We define static I/O mapping that stems only from the channel pathways such as displays. %. Such interactions would include displays as seen before.
% \subsubsection{Dynamic Interactivity}
Dynamic I/O mapping is expected when dynamic user interaction with the object or input light is present such as plugging in the switch knob, or blocking the channel by grasping a target object. Here, the end user has more control over the light that can be input and the dynamic re-routing of channels.
Various mechanisms can enable using light motion together with dynamic re-routing. For example, in the slider mechanism (see Figure \ref{fig:slider}), moving the slider knob causes movement of output light in relation to the input light at the bottom of the slider. A user may block light from a channel to sense which channel is being blocked or transmitting light based on channel parameters, or even  
% For example, a many-to-one mapping with channels ending on different faces of an object can be used to detect touch or blocking. A similar example of object orientation detection can be seen in Section \ref{sec:apps}.
decouple an existing channel for dynamic interactivity (Figure \ref{fig:bikehandle}).

%\hd{aligning the surface channel, depending on the surface of the object chosen as input or output}
%\section{Applications}
\label{sec:apps}

Modularity and flexibility enable a wide range of interactive and functional applications. By embedding optical channels within everyday objects, we can translate physical interactions such as squeezing, rotating, or changing orientation into meaningful visual outputs using light. %In this section, we demonstrate how these principles are applied in various applications, including bicycle handle signals, lane assistance indicators, and orientation sensors.


 % may also be sensed









%\subsection{Bistable switch/push button}




%\subsection{Image manipulation}

\section{Discussion \& Future Work}

% In this section, we reflect on the limitations of our current fabrication approach and propose directions for future development.% These include improving modular assembly through interlocking weave structures, discussing our method for single-extruder printers, and addressing geometric and material constraints that impact light transmission in FDM-printed optical channels.

\subsection{Weave Design for Assembly}

Currently, we assemble the partitions using external adhesive, however, the connections are not seamless.
To address the challenge posed by the lack of physical connections between the partitions, for future work we will encapsulate a weave channel design within the workflow. This approach involves restructuring the object into interlocking weave patterns that provide spaces for the optical channels. % but also significantly enhance the structural integrity of the printed object.
The weave design is based on volumetric Bravais weaves defined by Yildiz et. al. \cite{yildiz_volweaves2025}, which are a generalization of Grunbaum and Shephard's matrix encodings of planar weaves \cite{grunbaum1988isonemal}. We utilize the simplest form of weaves which is the plain weave due to its suitability for creating a regular, interlocking structure. To construct a plain weave, we instance two line segments that are perpendicular to each other on a 2D grid. Then we connect the neighboring weave elements to form a plain weave. A $3\times3$ plain weave structure is used in Figure \ref{fig:weavedicussion} to create light channel partitions.

There are other type of weaves such as satins and twill. All the periodic weaves that hang together up to $20\times20$ repeating domain \cite{grunbaum1988catalogue} have been cataloged. This flexibility enables the creation of custom weaves tailored to specific design requirements. Importantly, because the curves used in the weave channels lie entirely within a plane, i.e. they lack out-of-plane tangents, the design process for these weaves is streamlined, making them easier to implement compared to freeform curves.
The weave design can accommodate optical channels and address the assembly related issues of the current partitioning method.

\begin{figure}[h]
    \centering
    \begin{tikzpicture}[node distance=2.5mm]
    \node (a1) [align=center, on grid]{\includegraphics[width=0.33\linewidth]{figures/weavediscussion/weaveinputs.jpg}};
    \node (a2) [align=center, on grid, right=25mm of a1]{\includegraphics[width=0.33\linewidth]{figures/weavediscussion/result.jpg}};
    \node (a2) [align=center, on grid, right=25mm of a2]{\includegraphics[width=0.33\linewidth]{figures/weavediscussion/result_layout.jpg}};
    \end{tikzpicture}
    \vspace{-2mm}
    \caption{Light channel partitioning with woven channels.}
    \Description{A visual of light channel partitioning using a weave-based structure. A Stanford bunny is segmented into multiple interlocking regions formed by planar weaves, where horizontal and vertical bands alternate, creating slots for optical channels to pass through. The first part of the image shows the placement of the optical channels in the Stanford bunny. The second part of the image shows the partitions made in an assembled Stanford bunny. Finally, the six partitions are shown separately with the optical channels in a flat orientation.}
    \label{fig:weavedicussion}
    \vspace{-3mm}
\end{figure}

%Despite these improvements, we encounter challenges related to unintended holes forming in the design. These defects arise from the thin walls of the weave elements, which are misinterpreted by the slicer software as gaps during the slicing process. This highlights the need for further refinement in wall thickness and slicer configuration to prevent such issues in future implementations. Another shortcoming of using weaves as a light channel is the assembly process requires flexible material to weave the partitions into each other.

Another approach to assemble the partitions is making use of smaller connector pieces distributed computationally across the partitions. % For this approach we first wanted to get explicit corner information for each partition where we can place the connector pins.
In order to get explicit locations for placing connectors, we tried to use the channels with very few points such that Voronoi tesselation of the points creates a coarse mesh. Using this coarse mesh, we can query corners of the cells and assign connectors to each cell.
This approach, however, presents challenges due to the low-resolution sampling of the curves, which leads to inconsistent neighboring information between curves. The ongoing exploration in creating parametric connectors would also make the assembly process more seamless.

% \begin{figure}
%     \centering
%     \includegraphics[width=0.4\linewidth]{figures/single_extruder.png}
%     \vspace{-1mm}
%     \caption{Single extruder channel embedment requires first printing a groove of the channel and then overextruding transparent material in that groove at the last layer.}
%     \Description{Diagram illustrating the single-extruder method  A groove is printed into the shell of an object, and the extruder shows transparent material being extruded. The over-extrusion of this transparent filament fills the preformed groove, creating an embedded optical channel.}
    
%     \label{fig:single_extruder}
%     \vspace{-2mm}
% \end{figure}


% \subsection{Printing with a Single Extruder Printer}

% Our work relies on multi-extruder printing for in-place printing of channels. While multi-extruder printers are becoming easily accessible, we develop a script for using a single extruder printer where material for the shell can be swapped out with transparent filament by pausing the printing process. This technique uses the G-code of the partition to be printed, and analyses it to find the location of the channel using G-code comments. Once the location and the path of the channel is known, the script adjusts the printing process by omitting the material where the channel should be, for every layer except the last one.
% Here, a pause command is added to the G-code to allow the user to swap out shell material with transparent filament. During the final layer (Figure \ref{fig:single_extruder}), the script increases the material extrusion output which fills the previously created groove with the transparent material, forming the channel.
% The channel then behaves like a single channel line and does not have control parameters such as layer height or line width. Future evaluations will consider channels created by this script in considering the effect of thickness, curvature, and length on light transmission.

\subsection{Geometry, Material, \& FDM Limitations}

As seen from the evaluations, channels printed in FDM present some geometric and material limitations. While we test the maximum length of 100mm in the length evaluation similar to prior work \cite{printed_optics_2012}, 3D printed channels do not transmit light over large distances as optical fibers do. 
%Further, inserting printed channels has better transmissibility than printed in-place channels (Appendix \ref{app:a}). This may be happening due to frequent material switches and faster cooling of the transparent material during that time, or even impurities in the transparent material through inconsistent purging of shell material. We plan to further evaluate the reason behind this transmission mismatch between in-place printed and inserted channels. % case\jk{do we guess why?}. 
%However, certain pathways cannot be inserted into, and certain materials are not favorable for insertion. For example, rigid PLA cannot be inserted into a curved pathway. To enable use of different materials and different pathways, printing in-place is needed.
Future evaluations will also look at the effect of multiple curves within the same pathway on light transmission, and the effect of printing temperature on different transparent materials. While printing, we found that PETG loses its transparency and becomes milky white at certain temperatures, printing speeds, and cooling speeds. Evaluating this effect on transparency and milkyness of material will enable more control parameters that can be utilized to transmit or block light based on material temperature and similar parameters.

Although our evaluations primarily used Giantarm PLA and Overture TPU, future work will extend to other transparent materials. Leveraging a luminance meter (lux meter) \cite{gandhi20173d}, we aim to systematically characterize additional filaments to expand the material options for optical channel fabrication.


%\subsection{Resolution of Display}
FDM brings in limitations due to its layer-by-layer printing approach. As discussed earlier, channels need to be printed as continuous strokes. Hence, we cannot print channels that have 3D curves such as a helix.
Lastly, the resolution of printing is limited due to FDM compared to other approaches such as PolyJet printing.
While, theoretically, we can print channels as thin as a printer's nozzle allows, we had trouble printing consistent channels thinner than 1.6mm due to changes in print bed alignment, and print speed causing channels to be pulled off the bed.
The minimum consistent thinness we have achieved using FDM is 1.6mm. Further evaluations for print speed, temperature, and alignment consistency are warranted to print consistent thinner optical channels.

%\jk{while theoretically we can 3D print the cores as fine as possible as the print nozzle's diameter allows, it is not possible because..}
\section{Conclusion}
We showcase how FDM can be used to print in-place optical channels, enhancing the accessibility and customization of FDM optical interactive devices. 
We systematically evaluated the geometric and material characteristics of printed channels in PLA and TPU. 
By utilizing a Voronoi-based segmentation script, we have overcome challenges in embedding optical channels within FDM-printed structures, enabling interactive designs and functionality. We characterize the various design possibilities within a design space outlining static and dynamic  interactions. % through adjustable optical routing and channeling.


\begin{acks}
We would like to thank Abul Al Arabi for his help and suggestions with the circuit.
\end{acks}

%%
%% The next two lines define the bibliography style to be used, and
%% the bibliography file.
\bibliographystyle{ACM-Reference-Format}
\bibliography{references}


%%
%% If your work has an appendix, this is the place to put it.
\appendix

% \section{Inserted Channel Evaluation Data}
% \label{app:a}

% \begin{figure*}[h]

%     \centering
%     \includegraphics[width=\linewidth]{figures/PLA_insert_eval.pdf}
%     \caption{Light transmission through PLA printed channels which were inserted into the holder afterwards. Each subfigure shows the photodiode readings for red and blue light across varying parameters (a) length, (b) curvature, and (c) thickness. No significant differences were observed across parameters.}
%      \Description{} 
%     \label{fig:pla_insert}
%     \vspace{-2mm}
% \end{figure*}
% \begin{figure*}[h]

%     \centering
%     \includegraphics[width=\linewidth]{figures/PETG_insert_eval.pdf}
%     \caption{Light transmission through PETG printed channels which were inserted into the holder afterwards. Each subfigure shows the photodiode readings for red and blue light across varying parameters (a) length, (b) curvature, and (c) thickness. Inserted PETG channles also show no significant difference.}
%      \Description{} 
%     \label{fig:petg_insert}
%     \vspace{-2mm}
% \end{figure*}
% \begin{figure*}[h]

%     \centering
%     \includegraphics[width=\linewidth]{figures/TPU_insert_eval.pdf}
%     \caption{Light transmission through TPU printed channels which were inserted into the holder afterwards. Each subfigure shows the photodiode readings for red and blue light across varying parameters (a) length, (b) curvature, and (c) thickness. Inserted TPU channels show no significant difference.}
%      \Description{} 
%     \label{fig:tpu_insert}
%     \vspace{-2mm}
% \end{figure*}

% For PLA, PETG, and TPU, the inserted channels showed minimal loss, with photodiode readings remaining relatively constant near the maximum for channel length, curvature, and thickness as seen in Figures \ref{fig:pla_insert}, \ref{fig:petg_insert}, and \ref{fig:tpu_insert}.


\section{Circuit Diagram}
\label{app:ckt}
In our experiments, the photodiode was connected in a reverse-bias configuration with a load resistor to form a transimpedance circuit. The cathode of the photodiode was tied to the 5 V supply, and the anode was connected to the Arduino analog input (A0). A 1 M$\Omega$ resistor was placed from the analog input to ground, so that the photocurrent generated under illumination flowed through the resistor, developing a voltage that the Arduino’s ADC could measure.
%In our experiments, the photodiode was connected as a voltage-divider circuit. One terminal of the photodiode was tied to the supply, and the other terminal was connected to an analog input on the Arduino (A0). A load resistor (1 M$\Omega$) was placed from the analog input node to ground. This forms a divider: the photocurrent generated by the diode under illumination flows through the load resistor, developing a voltage that is sampled by the microcontroller’s ADC.
    
    \begin{figure}[h]
    \centering
\includegraphics[width=0.9\linewidth]{figures/circuit2.png}
    \caption{Diagram of circuit used for evaluations.} 
    \Description{the photodiode was connected in a reverse-bias configuration with a load resistor to form a transimpedance circuit. The cathode of the photodiode was tied to the 5 V supply, and the anode was connected to the Arduino analog input. A 1 M$\Omega$ resistor was placed from the analog input to ground, so that the photocurrent generated under illumination flowed through the resistor, developing a voltage that the Arduino’s ADC could measure.} 
    \label{fig:circuit}
    \vspace{-3mm}
\end{figure}


\section{Printing Parameters}
\label{app:params}
Table \ref{tab:printing_params} outlines the parameters utilized for printing length, thickness and curvature evaluation samples.

\begin{table}[h]
\centering
\caption{3D printing parameters used for PLA \& TPU samples for length, thickness, and curvature evaluations.}
\Description{Three columns: Parameter, PLA, and TPU. Each parameter and the corresponding values are written in the columns.}
\label{tab:printing_params}
\begin{tabular}{lcc}
\hline
\textbf{Parameter}            & \textbf{PLA}                 & \textbf{TPU} \\ \hline
Nozzle diameter               & 0.4 mm                       & 0.4 mm \\
Layer height                  & 0.4 mm                       & 0.4 mm \\
Line width                  & 0.41 mm                       & 0.41mm \\
First layer height            & 0.4 mm                      & 0.4 mm \\
Nozzle temperature            & 220\,$^{\circ}$C             & 228\,$^{\circ}$C \\
Bed temperature               & 55\,$^{\circ}$C              & 60\,$^{\circ}$C \\
Infill density (channel)                & 100\% (concentric) & 100\% (concentric ) \\
Infill density (part)                & 15\% (grid) & 50\% (lines ) \\
Cooling fan                   & 100\% after 1st layer        & 100\% after 1st layer \\
\hline
\end{tabular}
\end{table}



\begin{figure*}[ht]
    \centering
    \begin{tikzpicture}[node distance=2mm]
    \node (s11) {\includegraphics[width=0.30\linewidth]{figures/steps_screenshot/1.jpg}};
    \node (s12) [right=of s11] {\includegraphics[width=0.30\linewidth]{figures/steps_screenshot/2.jpg}};
    \node (s13) [right=of s12] {\includegraphics[width=0.30\linewidth]{figures/steps_screenshot/3.jpg}};
    \node (s21) [below=of s11] {\includegraphics[width=0.30\linewidth]{figures/steps_screenshot/4.jpg}};
    \node (s22) [right=of s21] {\includegraphics[width=0.30\linewidth]{figures/steps_screenshot/5.jpg}};
    \node (s23) [right=of s22] {\includegraphics[width=0.30\linewidth]{figures/steps_screenshot/6.jpg}};
    \node (s31) [below=of s21] {\includegraphics[width=0.30\linewidth]{figures/steps_screenshot/7.jpg}};
    \node (s32) [right=of s31] {\includegraphics[width=0.30\linewidth]{figures/steps_screenshot/8.jpg}};
    \node (s33) [right=of s32] {\includegraphics[width=0.30\linewidth]{figures/steps_screenshot/9.jpg}};
    \node (s41) [below=of s31] {\includegraphics[width=0.30\linewidth]{figures/steps_screenshot/10.jpg}};
    \node (s42) [right=of s41] {\includegraphics[width=0.30\linewidth]{figures/steps_screenshot/11.jpg}};
    \node (s43) [right=of s42] {\includegraphics[width=0.30\linewidth]{figures/steps_screenshot/12.jpg}};
    \node (l11) [draw=black,fill=white,text=black, circle, below right=0pt of s11.north west]{1};
    \node (l12) [draw=black,fill=white,text=black, circle, below right=0pt of s12.north west]{2};
    \node (l13) [draw=black,fill=white,text=black, circle, below right=0pt of s13.north west]{3};
    \node (l21) [draw=black,fill=white,text=black, circle, below right=0pt of s21.north west]{4};
    \node (l22) [draw=black,fill=white,text=black, circle, below right=0pt of s22.north west]{5};
    \node (l23) [draw=black,fill=white,text=black, circle, below right=0pt of s23.north west]{6};
    \node (l31) [draw=black,fill=white,text=black, circle, below right=0pt of s31.north west]{7};
    \node (l32) [draw=black,fill=white,text=black, circle, below right=0pt of s32.north west]{8};
    \node (l33) [draw=black,fill=white,text=black, circle, below right=0pt of s33.north west]{9};
    \node (l41) [draw=black,fill=white,text=black, circle, below right=0pt of s41.north west]{10};
    \node (l42) [draw=black,fill=white,text=black, circle, below right=0pt of s42.north west]{11};
    \node (l43) [draw=black,fill=white,text=black, circle, below right=0pt of s43.north west]{12};
    \end{tikzpicture}
    \caption{Screenshots from User Interface in Blender}
    \Description{Image shows the UI provided by the Blender and the implemented Add-On Interface. It shows 12 steps that goes through how to design channels and prepare the object for printing. Steps 1-6 shows how to designs first channel, steps 7-9 shows how to design the second channel, steps 10-11 shows how to apply partitioning and finally step 12 shows how to reorient the partitions for printing.}
    \label{fig:blenderui}
\end{figure*}



  \section{User Interface for Workflow Inside Blender\label{app:blender}}
As shown in Figure \ref{fig:blenderui} all of the design steps are implemented in Blender. Here is the summary of the steps shown in the figure:

\begin{enumerate}[noitemsep, leftmargin=*]
    \item Select object of interest in Blender viewport
    \item From the Add-On Menu, go into the \textit{Curve Partitioning Tool} and click on \verb|Add Canvas and Draw Channels|. Then select two points on the object. After selecting two points, tool will generate a plane that can be rotated with the mouse movement.
    \item Once satisfied with the plane orientation, confirm the orientation by clicking right mouse button. 
    \item Context is set to \verb|curve drawing mode| after confirming plane orientation
    \item Curves are drawn on the plane, making it easier for users to draw planar curves in 3D. At this stage all curve tools are available, whether it be deleting, drawing multiple curves, or changing vertex locations. %, all are permitted.
    \item Once satisfied with the drawn curves, they confirm the curve by pressing \verb|escape|
    \item To draw another curve, go over the same actions with clicking on \verb|Add Canvas and Draw Channels| on the Add-On menu, select points and orient the drawing plane.
    \item Draw a new curve
    \item Once satisfied with all the curves finalize the curves.
    \item Click on \verb|Apply Partitioning| on the Add-On menu to apply curve based partitioning to the object using Voronoi.
    \item Last step is to apply a procedural modifier to the object that partitions the object. There are several parameters that can be modified such as channel size, gap between the partitions, etc.
    \item Click on \verb|Apply 2D Printing Layout| on the Add-On menu to align each of the partitions and adjust the spacing between the partitions. Since all the workflow is implemented in a procedural geometry creation way, user can change parameters by going into the modifier and changing the values.
\end{enumerate}
\end{document}
\endinput
%%
%% End of file `sample-sigconf-authordraft.tex'.
